\documentclass[11pt,a4paper]{article}

% ── Compilar con: tectonic algoritmos.tex  (motor XeTeX, UTF-8 nativo) ──
\usepackage{fontspec}
\usepackage{amsmath,amssymb}
\usepackage[margin=2.3cm]{geometry}
\usepackage{xcolor}
\usepackage{booktabs}
\usepackage{tabularx}
\usepackage{longtable}
\usepackage{listings}
\usepackage{algorithm}
\usepackage{algpseudocode}
\usepackage{enumitem}
\usepackage{underscore}
\usepackage{titlesec}
\usepackage[hidelinks]{hyperref}

% ── Colores ──
\definecolor{cazul}{HTML}{2C5F8A}
\definecolor{cverde}{HTML}{2E8B57}
\definecolor{cnaranja}{HTML}{C1660A}
\definecolor{cgris}{HTML}{5B6770}
\definecolor{cfondo}{HTML}{EDF2F7}
\definecolor{ccodefondo}{HTML}{F6F8FA}

% ── Numeración profunda: 1 / 1.1 / 1.1.1 / 1.1.1.1 ──
\setcounter{secnumdepth}{3}
\setcounter{tocdepth}{3}
\titleformat{\paragraph}[hang]{\normalfont\normalsize\bfseries\color{cazul}}{}{0pt}{}
\titlespacing*{\paragraph}{0pt}{2.2ex plus 1ex minus .2ex}{1ex plus .2ex}

\renewcommand{\contentsname}{Índice}
\renewcommand{\tablename}{Tabla}

% ── Código ──
\lstset{
  basicstyle=\ttfamily\footnotesize,
  backgroundcolor=\color{ccodefondo},
  frame=single,
  rulecolor=\color{cgris!40},
  breaklines=true,
  columns=fullflexible,
  keepspaces=true,
  showstringspaces=false,
  xleftmargin=6pt,
  framexleftmargin=4pt,
}
\newcommand{\code}[1]{\texttt{\small #1}}

% ── Pseudocódigo en español ──
\algrenewcommand\algorithmicrequire{\textbf{Entrada:}}
\algrenewcommand\algorithmicensure{\textbf{Salida:}}
\algrenewcommand\algorithmicend{\textbf{fin}}
\algrenewcommand\algorithmicif{\textbf{si}}
\algrenewcommand\algorithmicthen{\textbf{entonces}}
\algrenewcommand\algorithmicelse{\textbf{si no}}
\algrenewcommand\algorithmicfor{\textbf{para}}
\algrenewcommand\algorithmicforall{\textbf{para cada}}
\algrenewcommand\algorithmicdo{\textbf{hacer}}
\algrenewcommand\algorithmicwhile{\textbf{mientras}}
\algrenewcommand\algorithmicreturn{\textbf{retornar}}
\algrenewcommand\algorithmicfunction{\textbf{función}}
\algrenewcommand\algorithmicprocedure{\textbf{proc}}
\algrenewcommand\algorithmicrepeat{\textbf{repetir}}
\algrenewcommand\algorithmicuntil{\textbf{hasta que}}
\algrenewcommand{\algorithmiccomment}[1]{\hfill\textcolor{cgris}{\textit{// #1}}}
\algnewcommand\algorithmicomit{\textbf{omitir vuelta}}
\algnewcommand\Omitir{\algorithmicomit{}}

% ── Bloques de pseudocódigo numerados y divisibles entre páginas ──
\newcounter{pcode}[section]
\renewcommand{\thepcode}{\thesection.\arabic{pcode}}
\newenvironment{pseudo}[1]{%
  \par\medskip\noindent
  \refstepcounter{pcode}%
  \colorbox{cfondo}{\parbox{\dimexpr\linewidth-2\fboxsep\relax}{%
    \small\textbf{\textcolor{cazul}{Pseudocódigo \thepcode.}} #1}}%
  \par\nobreak\smallskip\begin{algorithmic}[1]%
}{\end{algorithmic}\par\medskip}

% ── Caja de nota ──
\newenvironment{nota}{%
  \par\medskip\noindent
  \begin{tabular}{@{}p{0.4em}@{\hspace{8pt}}p{\dimexpr\linewidth-2em\relax}@{}}
  \textcolor{cnaranja}{\rule{3pt}{1.1em}} &
}{\end{tabular}\par\medskip}

\newcommand{\bigsep}{\par\vspace{3pt}}

\title{\vspace{-1.4cm}\textbf{Algoritmos de Alginvesting}\\[6pt]
  \large X4 — Backtester \quad$\cdot$\quad X5 — Cerebro macro\\[4pt]
  \normalsize\color{cgris} Pseudocódigos comentados, de lo global a lo particular}
\author{Proyecto Alginvesting}
\date{\today}

\begin{document}
\maketitle
\thispagestyle{empty}

\begin{center}
\begin{minipage}{0.86\linewidth}
\small\color{cgris}
Este documento describe con precisión operativa los dos módulos más grandes y menos
visibles del sistema: \code{X4_backtester.py} (simulador histórico de la estrategia) y
\code{X5_macro_brain.py} (modelo que aprende qué parámetros conviene usar en cada
contexto). Cada sección parte de la idea global y baja al detalle línea a línea.
Los conceptos técnicos se definen en una nota al pie la primera vez que aparecen,
para no obligar a saltar al glosario.
\end{minipage}
\end{center}

\vspace{8pt}
\tableofcontents
\newpage

% ═══════════════════════════════════════════════════════════════════════════
\section{Panorama general}
% ═══════════════════════════════════════════════════════════════════════════

\subsection{Qué problema resuelve el sistema completo}

Alginvesting compra en niveles de precio que el mercado ya respetó antes. La idea de
fondo es simple: si un activo bajó muchas veces hasta cierto precio y desde ahí rebotó,
ese precio concentra órdenes de compra reales; volver a comprar ahí tiene, en promedio,
mejor relación entre ganancia y riesgo que comprar en cualquier punto al azar. El
sistema mantiene decenas o cientos de esos niveles activos por activo, coloca órdenes de
compra pendientes en ellos y las administra con un mecanismo de protección de
ganancias.

\bigsep
Ese objetivo se reparte en seis scripts. La Tabla~\ref{tab:pipeline} los sitúa. Este
documento cubre solo los dos últimos, pero es imposible entenderlos sin saber qué le
entrega cada uno de los anteriores.

\begin{table}[h]
\centering
\small
\begin{tabularx}{\linewidth}{@{}l l X@{}}
\toprule
\textbf{Script} & \textbf{Rol} & \textbf{Qué produce} \\
\midrule
X0 & Datos y soportes & Descarga velas\footnote{\textbf{Vela} (\textit{candle}): resumen
de todo lo que pasó con el precio en un intervalo de tiempo fijo. Se describe con cuatro
números —apertura (\textit{Open}), máximo (\textit{High}), mínimo (\textit{Low}) y cierre
(\textit{Close})— más el volumen negociado. El conjunto se abrevia OHLCV. \textbf{H1} =
velas de una hora; \textbf{M1} = velas de un minuto.} horarias y calcula, por optimización,
los $N$ soportes óptimos de cada activo. \\
X1 & Trading en vivo & Lee los soportes, coloca y gestiona órdenes reales en MetaTrader~5. \\
X2 & Score fundamental & Un número en $[0,1]$ por activo que resume su salud fundamental. \\
X3 & Features técnicas & Indicadores técnicos por vela (medias móviles, RSI, ATR, etc.). \\
\textbf{X4} & \textbf{Backtester} & \textbf{Simula X0+X1 sobre el pasado; produce trades, eventos y curva de capital.} \\
\textbf{X5} & \textbf{Cerebro macro} & \textbf{Aprende qué parámetros usar según el contexto y los escribe en un archivo.} \\
\bottomrule
\end{tabularx}
\caption{Los seis scripts del sistema. Este documento cubre X4 y X5.}
\label{tab:pipeline}
\end{table}

\subsection{Por qué X4 y X5 se explican juntos}

X4 y X5 se escribieron como dos programas separados, pero funcionan como uno solo. La
relación es circular y conviene fijarla desde ya:

\begin{itemize}[leftmargin=1.4em,itemsep=2pt]
\item \textbf{X5 necesita datos que no existen.} Para aprender qué parámetros funcionan
mejor en cada contexto hacen falta miles de operaciones cerradas, cada una etiquetada con
el contexto de mercado y los parámetros vigentes cuando se abrió. En vivo eso tomaría
años. X5 los fabrica corriendo X4 sobre el histórico.
\item \textbf{X4 necesita saber con qué parámetros simular.} Si siempre simulara con los
mismos, el conjunto de datos resultante no tendría variedad y ningún modelo podría
aprender de él. Por eso, en modo \code{--x5}, X4 cambia de parámetros cada pocos días
simulados: a veces al azar, a veces preguntándole al modelo que X5 ya entrenó.
\end{itemize}

\bigsep
En una frase: \textbf{X4 es el laboratorio y X5 es el experimentador}. X5 diseña las
condiciones del experimento, X4 lo ejecuta y devuelve los resultados, X5 aprende de ellos
y diseña el siguiente. La Sección~\ref{sec:interaccion} desarma ese bucle paso a paso, una
vez que ambos módulos se hayan explicado por separado.

\subsection{Cómo leer este documento}

La estructura es deliberadamente descendente. Cada uno de los dos capítulos centrales
(Secciones~\ref{sec:x4} y~\ref{sec:x5}) sigue el mismo orden:

\begin{enumerate}[leftmargin=1.6em,itemsep=2pt]
\item \textbf{Propósito y contrato}: qué entra, qué sale, qué garantiza y qué no.
\item \textbf{Arquitectura}: las piezas internas y cómo se relacionan.
\item \textbf{Pseudocódigo global}: el algoritmo completo en una página, sin detalles.
\item \textbf{Detalle}: cada pieza del pseudocódigo global, expandida.
\item \textbf{Parámetros}: qué controla cada perilla y qué pasa al moverla.
\end{enumerate}

\bigsep
Los pseudocódigos están numerados por sección (Pseudocódigo 2.1, 2.2, \dots) y nombran las
funciones reales del archivo \code{.py} correspondiente, de modo que se pueda saltar del
documento al código sin traducción intermedia. Cuando una línea del pseudocódigo simplifica
algo relevante, hay una nota al pie o un párrafo inmediatamente después.

\subsection{Vocabulario del dominio}
\label{sec:vocab}

Todo el documento usa las siglas de la Tabla~\ref{tab:vocab}. Están definidas también en
\code{CLAUDE.md}, pero se repiten aquí porque el resto del texto las da por sabidas.

\begin{table}[h]
\centering
\small
\begin{tabularx}{\linewidth}{@{}l X@{}}
\toprule
\textbf{Sigla} & \textbf{Significado} \\
\midrule
Soporte & Nivel de precio en el que se coloca una orden de compra pendiente. \\
$N$ & Cantidad de soportes activos por activo. \\
O0 & Orden potencial: su soporte todavía está demasiado cerca del precio actual, así que aún no se coloca. \\
OE & \textbf{Orden en Espera}: orden de compra límite\footnote{\textbf{Orden límite de compra}
(\textit{buy limit}): instrucción al bróker de comprar automáticamente si el precio baja
hasta un nivel indicado. Mientras el precio no llegue, la orden existe pero no consume
capital: solo reserva margen cuando se ejecuta.} ya colocada, esperando que el precio baje hasta ella. \\
OA & \textbf{Orden Abierta}: la OE se ejecutó; hay una posición viva ganando o perdiendo. \\
OC & \textbf{Orden Cerrada}: la posición se cerró (por protección de ganancias, por pérdida máxima o por stop). \\
$A$ & Ganancia flotante en dólares que debe alcanzar una posición para que se active por primera vez su stop protector. \\
$B$ & Distancia en dólares a la que el stop protector sigue al precio una vez activado. \\
$L$ & Tamaño efectivo de la posición: lote $\times$ unidades por lote. Convierte diferencias de precio en dólares. \\
FO & \textbf{Función Objetivo}: número que mide qué tan bueno es un conjunto de $N$ soportes. X0 la maximiza. \\
$T$ & Ventana de días de historia usada para calcular soportes. \\
\bottomrule
\end{tabularx}
\caption{Vocabulario mínimo. La progresión natural de una orden es O0 $\to$ OE $\to$ OA $\to$ OC.}
\label{tab:vocab}
\end{table}

\subsubsection{La mecánica de una operación, con números}

Antes de entrar en el código conviene tener en la cabeza una operación completa, porque
todo X4 no es más que repetir esto millones de veces.

\bigsep
Supón NVDA cotizando a \$180. Hay un soporte en \$174. Los parámetros son $A=6$, $B=2$,
lote $=0{,}01$, unidades por lote $=100$, es decir $L = 0{,}01 \times 100 = 1$ (cada dólar
de movimiento del precio vale un dólar en la cuenta).

\begin{enumerate}[leftmargin=1.6em,itemsep=3pt]
\item \textbf{Creación (O0 $\to$ OE).} La distancia entre el precio y el soporte, medida en
dólares, es $(180 - 174) \times L = 6$. Como $6 \ge A = 6$, la orden se coloca. Si el
soporte estuviera en \$176, la distancia sería \$4 y la orden quedaría en O0: demasiado
cerca, no vale la pena inmovilizar una orden ahí.
\item \textbf{Ejecución (OE $\to$ OA).} El precio baja y en algún momento una vela toca
\$174. La orden se ejecuta: ahora hay una posición comprada a \$174.
\item \textbf{Activación del stop protector (paso C).} El precio sube a \$181. La ganancia
flotante es $(181-174)\times L = 7 \ge A = 6$, así que se activa por primera vez un stop
en $181 - B/L = 181 - 2 = \$179$. Desde este momento, la operación no puede terminar en
pérdida.
\item \textbf{Seguimiento.} El precio sigue a \$184: el stop se mueve a $184-2 = \$182$.
Nunca baja; si el precio retrocede a \$183, el stop se queda en \$182.
\item \textbf{Cierre (OA $\to$ OC).} El precio cae a \$182 y toca el stop. La posición se
cierra con $(182-174)\times L = \$8$ de ganancia. Se registra una fila OC.
\end{enumerate}

\bigsep
El camino alternativo es el malo: el precio nunca sube lo suficiente para activar el stop y
en cambio cae hasta que $(174 - \text{precio}) \times L > \code{PERDIDA_MAX}$. Entonces la
posición se cierra a pérdida (paso D). Toda la estrategia vive en la tensión entre esos dos
finales, y todos los parámetros que X5 optimiza mueven la frontera entre uno y otro.

\newpage
% ═══════════════════════════════════════════════════════════════════════════
\section{X4 — Backtester}
\label{sec:x4}
% ═══════════════════════════════════════════════════════════════════════════

\subsection{Propósito y contrato}

\subsubsection{Qué hace, en una frase}

X4 toma el historial de precios ya descargado, se para en la primera vela y avanza hora a
hora hasta el presente, comportándose exactamente como se habría comportado el sistema
real en cada uno de esos momentos: calcula soportes con la información disponible hasta
esa hora, coloca órdenes, las ejecuta cuando el precio las alcanza, mueve stops, cierra
posiciones y anota todo. Al final entrega un historial completo de operaciones y una curva
de capital.

\bigsep
Esto es \textit{backtesting}\footnote{\textbf{Backtesting}: evaluar una estrategia
aplicándola sobre datos históricos como si se estuviera operando en ese momento. Su
validez depende por completo de no usar información del futuro (ver
Sección~\ref{sec:x4-nolookahead}).}. Su utilidad es doble: permite estimar si la estrategia
gana o pierde dinero antes de arriesgar capital real, y —el uso dominante hoy— genera el
conjunto de datos de entrenamiento de X5.

\subsubsection{Qué garantiza: ausencia de información del futuro}
\label{sec:x4-nolookahead}

El error clásico de un backtester es dejar entrar, sin darse cuenta, información que en el
momento simulado todavía no existía. X4 se protege de eso en tres puntos concretos:

\begin{itemize}[leftmargin=1.4em,itemsep=2pt]
\item \textbf{Soportes.} Cuando en el instante simulado $t$ hay que recalcular soportes, se
llama a la misma función de X0 pasándole \code{fecha_hora_max = t}. Esa función recorta el
CSV a $\{$filas con fecha $\le t\}$ antes de optimizar. Un soporte nunca se apoya en una
vela que aún no ocurrió.
\item \textbf{Features de contexto.} El fotograma de indicadores técnicos (X3) y de score
fundamental (X2) que acompaña a cada operación se calcula también con el corte
\code{df.index <= ts}, y para X2 se elige la entrada del historial cuya fecha sea $\le$ la
fecha simulada.
\item \textbf{Columnas prohibidas en el modelo.} X5 excluye explícitamente del vector de
entrada cualquier columna terminada en \code{_oc} (información del momento del cierre), que
por definición no está disponible cuando hay que decidir.
\end{itemize}

\subsubsection{Qué no simula: los supuestos que hay que tener presentes}

Un backtester es siempre una simplificación. Los supuestos de X4, ordenados de más a menos
relevante:

\begin{enumerate}[leftmargin=1.6em,itemsep=3pt]
\item \textbf{Sin deslizamiento ni comisiones.} Una orden límite en \$174 se ejecuta a
\$174 exactos, salvo el caso de hueco (Sección~\ref{sec:x4-pasoB}). No se descuentan
comisiones ni costo de financiamiento nocturno. El resultado simulado es, por tanto, un
techo optimista.
\item \textbf{Liquidez infinita.} Se asume que siempre hay contraparte para el tamaño
operado. Con los lotes mínimos que usa el sistema es un supuesto razonable.
\item \textbf{Intra-vela sintético.} Dentro de una vela de una hora, el orden real en que
el precio visitó su máximo y su mínimo es desconocido. X4 lo reconstruye tomando prestado
un tramo real de sesenta minutos y reescalándolo (Sección~\ref{sec:x4-intravela}). Es una
aproximación, y es la mayor fuente de incertidumbre del simulador.
\item \textbf{Sin cierre forzado al final.} Cuando se acaban los datos, las posiciones
abiertas quedan abiertas: no se liquidan a precio de mercado. El capital final reportado es
saldo realizado, no valor de liquidación.
\end{enumerate}

\subsubsection{Los dos modos de ejecución}
\label{sec:x4-modos}

El mismo archivo se usa para dos cosas distintas, y confundirlas es la fuente habitual de
desorientación. La Tabla~\ref{tab:x4modos} las separa.

\begin{table}[h]
\centering
\small
\begin{tabularx}{\linewidth}{@{}l X X@{}}
\toprule
 & \textbf{Modo versión} & \textbf{Modo X5} \\
\midrule
Invocación & \code{python X4_backtester.py --version V1} & \code{python X4_backtester.py --x5 --activo BTCUSD} \\
Configuración & \code{resources/x4/versionV1/config_V1.py} & \code{scripts/config_x5_default.py} \\
Parámetros & Fijos durante toda la corrida & Se regeneran cada pocos días simulados \\
Pregunta que responde & ¿Cuánto habría rendido esta configuración exacta? & ¿Cómo rinden \textit{muchas} configuraciones en \textit{muchos} contextos? \\
Salida principal & \code{trades.json}, \code{equity_global.csv} & \code{resources/x5/\{ACTIVO\}_store.csv} \\
Activos & Todos los de la versión & Uno solo (aislamiento para paralelizar) \\
\bottomrule
\end{tabularx}
\caption{Los dos modos de X4. El interruptor es el argumento \code{--x5}.}
\label{tab:x4modos}
\end{table}

\bigsep
El modo versión existe para responder preguntas puntuales y reproducibles: cada versión
congela su propia copia de la configuración en un archivo aparte, de modo que cambiar
\code{config.py} meses después no altere el resultado de una corrida antigua. El script
auxiliar \code{X4B_crear_version_backtesting.py} crea ese andamiaje (carpeta, plantilla de
configuración, registro de la versión en \code{config.py}) a partir de un nombre.

\subsubsection{Entradas y salidas}

\begin{table}[h]
\centering
\small
\begin{tabularx}{\linewidth}{@{}l l X@{}}
\toprule
\textbf{Dirección} & \textbf{Archivo} & \textbf{Contenido} \\
\midrule
Entrada & \code{Data/\{ACTIVO\}.csv} & Velas H1 con OHLCV. Obligatorio. \\
Entrada & \code{Data_minuto/\{ACTIVO\}.csv} & Velas M1. Opcional: sin él, la simulación intra-vela se degrada. \\
Entrada & \code{resources/x2/x2_history.json} & Historial de scores fundamentales (solo modo X5). \\
Salida & \code{checkpoint.json} & Estado completo para reanudar. \\
Salida & \code{trades.json} & Una entrada por operación cerrada. \\
Salida & \code{events.json} & Bitácora de todo: creación, cancelación, cambio de stop, cierre. \\
Salida & \code{equity_global.csv} & Una fila por vela: saldo, patrimonio, margen. \\
Salida & \code{equity_activos.csv} & Una fila por vela y activo: ganancia cerrada, flotante y total. \\
Salida & \code{conjuntos_N/\{V\}_\{N\}_bt.json} & Caché de soportes indexado por fecha. \\
Salida & \code{\{ACTIVO\}_store.csv} & Filas de entrenamiento para X5 (solo modo X5). \\
\bottomrule
\end{tabularx}
\caption{Contrato de archivos de X4. Todas las rutas salen de la configuración activa.}
\end{table}

\subsection{Arquitectura interna}

\subsubsection{El estado: una sola estructura que lo contiene todo}
\label{sec:x4-estado}

X4 no usa clases. Todo el estado de la simulación vive en un único diccionario anidado que
se pasa por referencia a cada función. Conocer su forma es conocer el programa:

\begin{lstlisting}
estado = {
  'capital': 3000.0,                 # saldo realizado en USD
  'ts_ultimo_procesado': '2026-03-04T11:00:00',
  'ts_ultimo_recalculo': '2026-03-01T09:00:00',
  'stop_out': False,
  'por_activo': {
     'BTCUSD': {
        'soportes': [61200.0, 61845.5, ...],       # lista ordenada de N precios
        'OE': {61200.0: {'lote':0.01, 'ts_creacion':...}, ...},
        'OA': {60310.0: {'lote':0.01, 'sl':60480.0, 'ts_apertura':...,
                         'ganancia_max':..., 'drawdown_max':...,
                         'duracion_velas':..., 'capital_apertura':...}, ...},
        'GC': 184.55,                              # ganancia cerrada acumulada
     },
     'ETHUSD': {...},
  },
}
\end{lstlisting}

\bigsep
Dos decisiones de diseño se leen directamente ahí:

\begin{itemize}[leftmargin=1.4em,itemsep=2pt]
\item \textbf{Las órdenes se indexan por precio, no por identificador.} \code{OE} y
\code{OA} son diccionarios cuya clave es el precio de la orden. Esto hace instantáneo
preguntar «¿ya hay una orden en este soporte?», que es la consulta más frecuente del bucle,
pero impone una restricción: no puede haber dos órdenes en el mismo precio exacto para el
mismo activo. En la práctica no ocurre, porque los soportes son distintos entre sí por
construcción.
\item \textbf{El estado es serializable.} Todo lo que hay ahí adentro son números, textos y
diccionarios: se puede volcar a JSON y recuperar tal cual. De eso vive el mecanismo de
reanudación.
\end{itemize}

\subsubsection{Checkpoint: parar y seguir}

Una pasada completa sobre varios años de velas horarias tarda horas. Cada 24 velas
procesadas —y siempre antes de un evento importante— X4 vuelca \code{estado} a
\code{checkpoint.json} y descarga a disco las listas de operaciones y eventos acumuladas en
memoria. Al arrancar, si encuentra un checkpoint y no se pidió \code{--reset}, lo carga y
salta todas las velas cuya fecha sea $\le$ \code{ts_ultimo_procesado}.

\begin{nota}
Un detalle fácil de pasar por alto: al saltar velas ya procesadas, X4 igual actualiza
\code{cierre_previo} de cada activo. Ese diccionario guarda el último precio conocido de
cada activo y se usa para valorar posiciones de activos que no tienen vela en el instante
actual (una acción de EE.UU.\ no cotiza de madrugada, pero su posición sigue existiendo).
Si no se actualizara durante el salto, el primer cálculo de patrimonio tras reanudar usaría
precios de hace meses.
\end{nota}

\bigsep
Las claves numéricas son el punto delicado de la serialización: JSON solo admite claves de
texto, así que al guardar se convierten a texto y al cargar se reconvierten con
\code{float(k)}. Esa conversión de ida y vuelta está en \code{_cargar_checkpoint} y
\code{_guardar_checkpoint} y debe mantenerse simétrica.

\subsubsection{Carga de datos}

Las velas H1 se cargan una vez, completas, en memoria: un \code{DataFrame} por activo,
ordenado por fecha, sin duplicados y recortado desde \code{fecha_inicio}. Se indexa por
fecha para que la consulta \code{df.loc[ts]} sea directa. Las velas M1 se cargan igual pero
sin indexar, porque no se consultan por fecha sino por posición aleatoria
(Sección~\ref{sec:x4-intravela}).

\bigsep
El eje temporal del bucle no es el de un activo, sino la \textbf{unión ordenada de todas
las fechas de todos los activos}. Así, un instante en que solo cotiza BTCUSD existe en el
bucle y se procesa solo para ese activo; el resto se valora al último precio conocido.

\subsection{El algoritmo global}

Con el estado y los datos ya presentados, el bucle completo cabe en una página. El
Pseudocódigo~\ref{alg:x4global} es el mapa: todo lo que viene después de esta subsección es
la expansión de alguna de sus líneas.

\begin{pseudo}{\code{ejecutar_backtest(cfg, reset, x5_mode)} — bucle principal}
\label{alg:x4global}
\Require configuración \code{cfg}, banderas \code{reset} y \code{x5_mode}
\Ensure archivos de resultados en disco y \code{estado} final
\State \textbf{datos_h1} $\gets$ cargar CSV horarios de todos los activos, recortados desde \code{fecha_inicio}
\State \textbf{datos_m1} $\gets$ cargar CSV de minutos (puede faltar por activo)
\If{\code{x5_mode}}
  \State decidir \code{reset} solo: pasada nueva si no hay checkpoint, si terminó, o si hubo quiebra
  \If{\code{reset}} borrar el estado adaptativo del optimizador de soportes \EndIf
\EndIf
\State \textbf{estado} $\gets$ checkpoint si \code{$\neg$ reset} y existe; si no, estado inicial con \code{capital_inicial}
\If{algún activo no tiene soportes}
  \If{\code{x5_mode}} elegir parámetros del primer tramo (Sec.~\ref{sec:x4-x5params}) \EndIf
  \State \textbf{recalcular_soportes}(estado, $t = $ \code{fecha_inicio}) \Comment{arranque en frío}
\EndIf
\State \textbf{todos_los_ts} $\gets$ unión ordenada de las fechas de todos los activos
\For{cada \textbf{ts} en \textbf{todos_los_ts}}
  \State actualizar \code{cierre_previo} de los activos con vela en \textbf{ts}
  \If{\textbf{ts} $\le$ \code{ts_ultimo_procesado}} \Omitir \Comment{ya simulada antes del checkpoint} \EndIf
  \If{han pasado \code{delta_recalculo_soportes} días desde el último recálculo}
     \If{\code{x5_mode}} elegir parámetros del nuevo tramo (explorar o explotar) \EndIf
     \State \textbf{recalcular_soportes}(estado, \textbf{ts})
     \State guardar avance y \Omitir \Comment{la vela del recálculo no se opera}
  \EndIf
  \State \textbf{precios} $\gets$ cierre de cada activo en \textbf{ts}, o su último cierre conocido
  \For{cada \textbf{activo} con vela en \textbf{ts}}
     \If{\code{x5_mode}} \textbf{ctx} $\gets$ fotograma de features X3 y X2 al instante \textbf{ts} \EndIf
     \State \textbf{procesar_vela}(\textbf{ts}, vela, \textbf{activo}, \dots) \Comment{pasos A--F}
     \If{\code{x5_mode} y toca registro periódico y hay posiciones abiertas}
        \State escribir fila \code{'periodico'} en el store de X5
     \EndIf
  \EndFor
  \State \textbf{mc} $\gets$ \textbf{registrar_equity}(\textbf{ts}, estado, \textbf{precios}) \Comment{paso G}
  \State \code{ts_ultimo_procesado} $\gets$ \textbf{ts}
  \If{patrimonio $\le 0$ \textbf{o} nivel de margen $\le$ \code{STOP_OUT_LEVEL}}
     \State registrar quiebra, guardar todo y \textbf{cortar}
  \EndIf
  \If{múltiplo de 24 velas} guardar checkpoint y volcar operaciones y eventos a disco \EndIf
\EndFor
\State guardar checkpoint final, volcar pendientes, imprimir resumen
\end{pseudo}

\bigsep
Tres observaciones sobre este esqueleto:

\paragraph{La vela del recálculo se congela.} Cuando toca recalcular soportes, X4 recalcula,
guarda y salta a la siguiente vela sin operar. El motivo es de coherencia interna: el
recálculo puede cancelar órdenes en espera y crear otras nuevas, y mezclar eso con la
ejecución de la misma vela dejaría un orden de eventos ambiguo (¿la orden que se canceló
alcanzó a ejecutarse?). Congelar la vela cuesta una hora de simulación cada varios días y
elimina la ambigüedad.

\paragraph{El disparo del recálculo es por umbral, no por hora exacta.} La condición es
«han pasado al menos $\delta$ días», evaluada en la primera vela disponible que la cumpla.
Una versión anterior exigía además una hora concreta del día, y eso hacía que activos sin
sesión de 24 horas —las acciones— pudieran no tener nunca una vela a esa hora y quedaran
con soportes congelados para siempre.

\paragraph{El orden interno de cada vela es fijo.} Primero se procesan todos los activos,
después se registra el patrimonio, después se evalúa la quiebra. Registrar el patrimonio
antes de operar daría una foto desfasada; evaluar la quiebra antes de registrar dejaría la
última fila del historial sin escribir.

\subsection{Los pasos A--F: la lógica de trading de una vela}

\code{_procesar_candle} es el corazón. Recibe una vela de un activo y aplica seis pasos en
orden estricto. Antes de ellos hace dos actualizaciones de contadores: incrementa la
duración en velas de cada posición abierta y actualiza su peor pérdida flotante usando el
mínimo de la vela.

\begin{pseudo}{\code{_procesar_candle(ts, vela, activo, \dots)}}
\Require vela con Open/High/Low/Close, estado del activo
\For{cada posición abierta \textbf{p}}
  \State \code{p.duracion_velas} $\gets$ \code{p.duracion_velas} $+\,1$
  \State \code{p.drawdown_max} $\gets \min($\code{p.drawdown_max}$, (\text{Low} - p_{ap}) \cdot L)$
\EndFor
\State \textbf{precio_saliente} $\gets$ \textbf{paso_A}(\dots) \Comment{cancelar OE huérfanas}
\State \textbf{usar_iv} $\gets$ \textbf{necesita_intravela}(\dots)
\If{\textbf{usar_iv} y hay datos de minutos}
  \State \textbf{simular_intravela}(\dots) \Comment{repite B--F sobre 60 velas de un minuto}
\Else
  \State \textbf{paso_B}(vela) \Comment{ejecutar OE alcanzadas}
  \State \textbf{paso_C}(High) \Comment{mover stops protectores}
  \State \textbf{paso_D}(Low) \Comment{cerrar por pérdida máxima}
  \State \textbf{paso_E}(Low) \Comment{cerrar por stop tocado}
  \State \textbf{paso_F}(Close) \Comment{crear OE nuevas}
\EndIf
\end{pseudo}

\subsubsection{Paso A — cancelar órdenes en espera huérfanas}

Cada recálculo produce un conjunto de soportes distinto. Una orden en espera colocada sobre
un soporte que ya no existe debe retirarse: el nivel dejó de ser interesante. El paso A
recorre las OE y elimina las que no coinciden con ningún soporte vigente.

\begin{pseudo}{\code{_paso_A} — cancelación de OE sin soporte}
\Ensure precio más alto entre las canceladas, o \textit{nada}
\State \textbf{a_eliminar} $\gets \{p \in$ OE $: p \notin$ soportes$\}$
\For{cada \textbf{p} en \textbf{a_eliminar}}
  \State registrar evento \code{OE_eliminada} con motivo \code{soporte_desactivado}
  \State borrar OE en \textbf{p}
\EndFor
\State \Return $\max($\textbf{a_eliminar}$)$ si no está vacío
\end{pseudo}

\bigsep
El valor de retorno no es decorativo. Si se acaba de retirar un buy limit que estaba
razonablemente cerca del precio, la zona inmediatamente inferior queda descubierta. El paso
F usa ese precio para permitir que un soporte cercano ocupe el hueco aunque no cumpla la
distancia mínima habitual (Sección~\ref{sec:x4-pasoF}). Es la misma regla de reemplazo que
aplica el sistema en vivo en \code{X1.crear_ordenes_espera}, y mantenerla idéntica en ambos
lados es lo que hace comparables los resultados simulados con los reales.

\subsubsection{Paso B — ejecutar las órdenes que el precio alcanzó}
\label{sec:x4-pasoB}

Una orden en espera a precio $p$ se ejecuta si el precio bajó hasta $p$. Con velas horarias
eso se traduce en $\text{Low} \le p$. Pero hay un caso especial que merece tratamiento
aparte: el hueco de mercado.

\paragraph{El caso normal.} Se recorren las OE con $p < \text{Open}$ de mayor a menor. Si
$\text{Low} \le p$, la orden se ejecuta al precio $p$ exacto.

\paragraph{El hueco.} Si una orden estaba en \$174 y el mercado abre en \$170 —viernes cierre
en \$176, lunes apertura en \$170 tras una mala noticia—, el precio nunca pasó por \$174: se
lo saltó\footnote{\textbf{Hueco} (\textit{gap}): salto de precio entre el cierre de una vela
y la apertura de la siguiente, sin cotizaciones intermedias. Típico de fines de semana y de
la apertura de mercados no continuos. Es la razón por la que una orden límite puede
ejecutarse a un precio peor que el pedido.}. En la realidad el bróker ejecuta todas las
órdenes límite pendientes por encima de la apertura al precio de apertura. X4 replica eso:
detecta el conjunto \code{oes_gap} $= \{p \in$ OE $: p \ge \text{Open}\}$ y ejecuta todas al
precio $\min($\code{oes_gap}$)$, no al precio pedido. El resultado es una entrada peor que la
planificada, que es exactamente lo que ocurre en vivo.

\begin{pseudo}{\code{_paso_B} — ejecución de OE, con y sin hueco}
\State \textbf{oes_gap} $\gets \{p \in$ OE $: p \ge \text{Open}\}$
\If{\textbf{oes_gap} no está vacío}
  \State \textbf{precio_gap} $\gets \min($\textbf{oes_gap}$)$
  \For{cada \textbf{p} en \textbf{oes_gap} ordenadas}
     \State \textbf{margen_nueva} $\gets$ \textbf{precio_gap} $\cdot$ lote $\cdot$ unidades $/$ apalancamiento
     \If{margen libre $-$ \textbf{margen_nueva} $<$ \code{MARGEN_LIBRE_MIN_BT}}
        \State cancelar la OE con motivo \code{margen_insuficiente}; \Omitir
     \EndIf
     \State mover OE $\to$ OA anotando \textbf{precio_gap} como precio de apertura
     \State registrar evento \code{OE_ejecutada} con \code{es_gap = verdadero}
  \EndFor
\EndIf
\For{cada \textbf{p} en OE con $p < \text{Open}$, de mayor a menor}
  \If{$\text{Low} \le p$}
     \State mismo control de margen; si no alcanza, cancelar y \Omitir
     \State mover OE $\to$ OA con precio de apertura $p$
     \State registrar evento \code{OE_ejecutada} con \code{es_gap = falso}
  \EndIf
\EndFor
\end{pseudo}

\paragraph{El control de margen.} Antes de convertir cada OE en OA se verifica que quede
margen libre\footnote{\textbf{Margen}: dinero que el bróker inmoviliza como garantía al
abrir una posición apalancada. Se calcula como valor de la posición dividido por el
apalancamiento. \textbf{Margen libre} = patrimonio $-$ margen usado: lo que queda disponible
para abrir posiciones nuevas o absorber pérdidas.} suficiente. Si no lo hay, la orden no se
ejecuta: se cancela y se registra el motivo. Esto reproduce el rechazo del bróker y evita
que la simulación abra posiciones imposibles.

\begin{nota}
El control se hace \textit{orden por orden}, recalculando el margen libre en cada
iteración, porque cada ejecución consume margen y cambia la respuesta para la siguiente.
En un hueco grande pueden ejecutarse decenas de órdenes en el mismo instante y la última no
tiene por qué caber.
\end{nota}

\subsubsection{Paso C — el stop protector móvil}

El mecanismo que convierte una ganancia flotante en ganancia asegurada. Tiene dos
regímenes:

\begin{itemize}[leftmargin=1.4em,itemsep=2pt]
\item \textbf{Sin stop todavía} (\code{sl = 0}). Se activa por primera vez solo si la
ganancia flotante calculada con el máximo de la vela alcanza $A$:
$(\text{High} - p_{ap}) \cdot L \ge A$. El stop se coloca en $\text{High} - B/L$.
\item \textbf{Con stop activo.} El nuevo candidato es $\text{High} - B/L$. Se adopta
\textit{solo si es mayor} que el actual. El stop sube con el precio y nunca baja.
\end{itemize}

\begin{pseudo}{\code{_paso_C} — trailing stop\protect\footnotemark}
\For{cada posición abierta en $p_{ap}$}
  \State \textbf{ga} $\gets (\text{High} - p_{ap}) \cdot L$ \Comment{ganancia flotante con el máximo de la vela}
  \State \code{ganancia_max} $\gets \max($\code{ganancia_max}$,$ \textbf{ga}$)$
  \If{\code{sl} $= 0$}
     \If{\textbf{ga} $\ge A$}
        \State \code{sl} $\gets \text{High} - B/L$
        \State reponer una OE en $p_{ap}$ si no existe \Comment{el nivel vuelve a quedar libre}
        \State registrar evento \code{SL_cambiado}
     \EndIf
  \Else
     \State \textbf{sl_nuevo} $\gets \text{High} - B/L$
     \If{\textbf{sl_nuevo} $>$ \code{sl}} \State \code{sl} $\gets$ \textbf{sl_nuevo}; registrar evento \EndIf
  \EndIf
\EndFor
\end{pseudo}
\footnotetext{\textbf{Trailing stop} (stop dinámico o «de arrastre»): orden de venta
automática cuyo precio de disparo sigue al precio de mercado mientras este sube, y se queda
quieta cuando baja. Su efecto es fijar un piso creciente para el resultado de la operación.}

\paragraph{La reposición de la orden en espera.} La línea sutil de este paso: cuando una
posición activa su primer stop, se vuelve a colocar una OE en su mismo precio de apertura.
La lectura es que ese nivel ya demostró funcionar —el precio bajó hasta ahí, rebotó y la
posición está en ganancia asegurada—, así que vale la pena tener otra orden esperando por si
el precio vuelve. Sin esta línea, cada soporte serviría para una sola operación por ciclo de
recálculo.

\subsubsection{Paso D — cierre por pérdida máxima}

Corte de pérdidas duro. Si la pérdida flotante medida con el mínimo de la vela supera el
umbral, la posición se cierra al mínimo:
\[
(p_{ap} - \text{Low}) \cdot L > \code{PERDIDA_MAX_BT} \;\Longrightarrow\; \text{cerrar a } \text{Low}
\]

\begin{pseudo}{\code{_paso_D} — cierre por pérdida máxima}
\State \textbf{a_cerrar} $\gets \{p_{ap} \in$ OA $: (p_{ap} - \text{Low}) \cdot L > \code{PERDIDA_MAX_BT}\}$
\For{cada $p_{ap}$ en \textbf{a_cerrar}}
  \State sacar la posición de OA; actualizar su peor pérdida flotante
  \State \textbf{retorno} $\gets (\text{Low} - p_{ap}) \cdot L$
  \State \code{capital} $\gets$ \code{capital} $+$ \textbf{retorno}; \quad \code{GC[activo]} $\gets$ \code{GC[activo]} $+$ \textbf{retorno}
  \State construir el registro de operación con motivo \code{perdida_max} y añadirlo a \code{trades}
  \If{\code{x5_mode}} construir la fila OC y escribirla en el store de X5 \EndIf
  \State registrar evento \code{posicion_cerrada}
\EndFor
\end{pseudo}

\bigsep
Nótese que el precio de cierre es el mínimo de la vela, no el nivel exacto donde se cruzó el
umbral. Es la convención conservadora: se asume el peor punto alcanzado dentro de la vela.
Cuando la simulación intra-vela está activa, el mismo paso se ejecuta sobre velas de un
minuto y el precio de cierre es mucho más ajustado.

\subsubsection{Paso E — cierre por stop tocado}

Idéntico en estructura al paso D, pero la condición es que el mínimo de la vela haya llegado
al stop, y el precio de cierre es el stop mismo (no el mínimo): una orden de stop se ejecuta
a su nivel.

\begin{pseudo}{\code{_paso_E} — cierre por stop protector}
\State \textbf{a_cerrar} $\gets \{p_{ap} \in$ OA $: \code{sl} > 0 \;\wedge\; \text{Low} \le \code{sl}\}$
\For{cada $p_{ap}$ en \textbf{a_cerrar}}
  \State \textbf{retorno} $\gets (\code{sl} - p_{ap}) \cdot L$ \Comment{normalmente positivo, por construcción del paso C}
  \State actualizar capital y ganancia cerrada del activo
  \State registrar la operación con motivo \code{trailing_stop} y, en modo X5, la fila del store
\EndFor
\end{pseudo}

\paragraph{Por qué D va antes que E.} Ambos miran el mínimo de la vela y ambos pueden
aplicar a la misma posición. Poner D primero significa que, si una posición cruza a la vez
su pérdida máxima y su stop, se registra como pérdida máxima. Con velas horarias es un
empate raro; con simulación intra-vela activa prácticamente desaparece, porque las
condiciones se evalúan minuto a minuto y una ocurre estrictamente antes que la otra.

\subsubsection{Paso F — crear órdenes en espera}
\label{sec:x4-pasoF}

Cierra el ciclo: recorre los soportes de mayor a menor y coloca una orden en espera en cada
uno que cumpla las condiciones.

\begin{pseudo}{\code{_paso_F} — creación de OE sobre soportes libres}
\Require precio de referencia $p_{ref}$ (cierre de la vela), \textbf{precio_saliente} del paso A
\State \textbf{margen_comprometido} $\gets 0$
\If{\textbf{precio_saliente} existe y \textbf{precio_saliente} $< p_{ref}$}
   \State \textbf{umbral_reemplazo} $\gets (p_{ref} + \textbf{precio_saliente})/2$
\EndIf
\For{cada \textbf{soporte} en orden decreciente}
  \If{ya hay OA u OE en \textbf{soporte}} \Omitir \EndIf
  \State \textbf{distancia_ok} $\gets (p_{ref} - \textbf{soporte}) \cdot L \ge A$
  \State \textbf{reemplazo_ok} $\gets$ \textbf{umbral_reemplazo} existe $\wedge$ \textbf{soporte} $<$ \textbf{umbral_reemplazo}
  \If{ni \textbf{distancia_ok} ni \textbf{reemplazo_ok}} \Omitir \Comment{queda en estado O0} \EndIf
  \State \textbf{margen_nueva} $\gets$ \textbf{soporte} $\cdot$ lote $\cdot$ unidades $/$ apalancamiento
  \If{margen libre $-$ \textbf{margen_comprometido} $-$ \textbf{margen_nueva} $<$ \code{MARGEN_LIBRE_MIN_BT}} \Omitir \EndIf
  \State \textbf{margen_comprometido} $\gets$ \textbf{margen_comprometido} $+$ \textbf{margen_nueva}
  \State crear OE en \textbf{soporte}; en modo X5, adjuntar el fotograma de features X3 y X2
  \State registrar evento \code{OE_creada}
\EndFor
\end{pseudo}

\paragraph{La condición de distancia.} $(p_{ref} - \text{soporte}) \cdot L \ge A$ dice: solo
coloca la orden si, cuando se ejecute, tendrá espacio para ganar al menos $A$ antes de
volver al precio actual. Colocar órdenes pegadas al precio produce entradas que se ejecutan
enseguida y quedan atrapadas sin margen de maniobra.

\paragraph{La regla de reemplazo.} El punto medio entre el precio actual y el buy limit
recién cancelado define una zona donde un soporte puede entrar aunque esté más cerca de lo
normal. Recupera la cobertura que el paso A acaba de retirar, sin abrir la puerta a llenar
toda la zona cercana de órdenes.

\paragraph{Margen comprometido dentro del ciclo.} Como el paso F puede crear muchas órdenes
en una sola pasada y el margen de una orden en espera no está realmente reservado hasta que
se ejecuta, se lleva un acumulador local \code{margen_comprometido}. Sin él, el paso
colocaría órdenes por un total que la cuenta no puede sostener si todas se ejecutaran a la
vez.

\subsubsection{Por qué este orden y no otro}

El orden A--F no es arbitrario. Se puede justificar como una cadena de dependencias:

\begin{enumerate}[leftmargin=1.6em,itemsep=2pt]
\item \textbf{A antes que todo}, porque opera sobre información que acaba de cambiar (los
soportes) y podría estar dejando órdenes vivas sobre niveles ya inválidos.
\item \textbf{B antes que C, D y E}, porque una orden que se ejecuta en esta vela debe poder
gestionarse en esta misma vela: si el precio cae fuerte, una posición recién abierta puede
tener que cerrarse de inmediato por pérdida máxima.
\item \textbf{C antes que E}, porque el stop debe estar actualizado antes de comprobar si el
precio lo tocó. Al revés, se cerrarían posiciones con un stop viejo.
\item \textbf{F al final}, porque necesita el estado ya depurado: sabe qué niveles quedaron
ocupados tras las ejecuciones de B y liberados tras los cierres de D y E, y cuánto margen
queda realmente.
\end{enumerate}

\subsection{Simulación intra-vela}
\label{sec:x4-intravela}

\subsubsection{El problema}

Una vela horaria dice que el precio abrió en 180, tocó 184, bajó hasta 176 y cerró en 181.
No dice en qué orden. Y el orden importa muchísimo: si primero subió a 184 y luego bajó a
176, el stop protector se activó y la operación cerró en ganancia; si primero cayó a 176 y
luego subió, quizá se cerró antes por pérdida máxima. Con las mismas cuatro cifras, dos
resultados opuestos.

\subsubsection{La solución adoptada}

Cuando el desenlace de la vela depende del orden interno, X4 toma un tramo real de sesenta
minutos del historial M1 del mismo activo —elegido al azar—, lo deforma para que encaje
exactamente en la vela horaria y simula los pasos B--F sobre cada uno de esos sesenta
minutos.

\bigsep
El reescalado es una transformación afín que preserva la forma del tramo pero fuerza sus
extremos:
\[
x' \;=\; \text{Low}_{H1} \;+\; (x - \text{Low}_{M1}) \cdot
\frac{\text{High}_{H1} - \text{Low}_{H1}}{\text{High}_{M1} - \text{Low}_{M1}}
\]
y luego se fija a mano la apertura de la primera vela de un minuto y el cierre de la última
para que coincidan con los de la vela horaria. El resultado es un camino de precios
plausible —con la microestructura de un tramo real de ese mismo activo— que respeta los
cuatro valores conocidos.

\begin{pseudo}{\code{_necesita_intravela} — ¿hace falta el detalle de minutos?}
\State \textbf{rango} $\gets \text{High} - \text{Low}$
\For{cada posición abierta}
  \If{\code{sl} $= 0$}
     \If{\textbf{rango} $\cdot L \ge A$ \textbf{o} $(p_{ap} - \text{Low})\cdot L \ge \code{PERDIDA_MAX_BT}$} \Return \textbf{verdadero} \EndIf
  \Else
     \State \textbf{sl_nuevo} $\gets \text{High} - B/L$
     \If{\textbf{sl_nuevo} $>$ \code{sl} \textbf{y} $\text{Low} \le$ \textbf{sl_nuevo}} \Return \textbf{verdadero} \EndIf
  \EndIf
\EndFor
\For{cada orden en espera en $p$}
  \If{$\text{Low} \le p$ \textbf{y} (\textbf{rango} $\cdot L \ge A$ \textbf{o} $(p - \text{Low})\cdot L \ge \code{PERDIDA_MAX_BT}$)} \Return \textbf{verdadero} \EndIf
\EndFor
\State \Return \textbf{falso}
\end{pseudo}

\bigsep
La prueba busca exactamente las situaciones ambiguas: una vela lo bastante amplia como para
que quepan dentro tanto la activación del stop como su ejecución posterior, o una posición
que podría a la vez alcanzar su stop y su pérdida máxima. Cuando la vela es tranquila, no
hay ambigüedad y se resuelve con los cuatro valores horarios, que es mucho más rápido.

\begin{pseudo}{\code{_simular_intravela}}
\If{no hay datos M1 o hay menos de 60 velas} \Return \EndIf
\State $t \gets$ entero aleatorio en $[0,\; |M1| - 60]$
\State \textbf{bloque} $\gets$ reescalar($M1[t : t+60]$, vela horaria)
\For{cada \textbf{vela_m1} en \textbf{bloque}}
  \State \textbf{paso_B}(\textbf{vela_m1}); \quad \textbf{paso_C}(High$_{m1}$); \quad \textbf{paso_D}(Low$_{m1}$)
  \State \textbf{paso_E}(Low$_{m1}$); \quad \textbf{paso_F}(Close$_{m1}$)
\EndFor
\end{pseudo}

\subsubsection{Limitaciones honestas de este enfoque}

\begin{itemize}[leftmargin=1.4em,itemsep=2pt]
\item \textbf{Es aleatorio.} Dos corridas idénticas dan resultados distintos, porque el
tramo de minutos elegido cambia. No hay semilla fija. Para el uso principal —generar
variedad para X5— eso es más virtud que defecto, pero invalida la comparación exacta entre
dos corridas del modo versión.
\item \textbf{El tramo prestado no corresponde al momento simulado.} Se toma de cualquier
punto del historial de minutos, no del instante que se está simulando. Aporta una
microestructura realista promedio, no la real de esa hora.
\item \textbf{Todas las velas del bloque comparten la marca de tiempo horaria.} Las
operaciones abiertas y cerradas dentro del bloque quedan fechadas con la hora de la vela
horaria. La duración en minutos de una operación intra-vela no es recuperable desde los
registros.
\item \textbf{Sin datos M1 se degrada en silencio.} Si falta el archivo de minutos, X4
avisa por consola y resuelve la vela con los cuatro valores horarios.
\end{itemize}

\subsection{Contabilidad de la cuenta}

\subsubsection{Las cuatro magnitudes}

En cada vela, tras procesar todos los activos, \code{_calcular_estado_cuenta} recorre las
posiciones abiertas y produce:

\begin{align*}
\text{ganancia flotante} &= \textstyle\sum_{\text{OA}} (p_{actual} - p_{ap}) \cdot \text{lote} \cdot \text{unidades} \\
\text{patrimonio} &= \text{saldo} + \text{ganancia flotante} \\
\text{margen usado} &= \textstyle\sum_{\text{OA}} \frac{p_{ap} \cdot \text{lote} \cdot \text{unidades}}{\text{apalancamiento}} \\
\text{nivel de margen} &= \frac{\text{patrimonio}}{\text{margen usado}} \times 100
\end{align*}

\bigsep
El \textbf{saldo} solo cambia cuando una posición cierra; el \textbf{patrimonio} cambia con
cada tick. El \textbf{nivel de margen}\footnote{\textbf{Nivel de margen} (\textit{margin
level}): patrimonio dividido por margen usado, en porcentaje. Es el indicador que vigila el
bróker: por debajo de cierto porcentaje empieza a cerrar posiciones por su cuenta
(\textit{margin call} y luego \textit{stop out}). Un nivel de 1000\,\% es holgado; 100\,\%
significa que el patrimonio apenas cubre la garantía exigida.} es la métrica de
supervivencia.

\subsubsection{Quiebra de la cuenta}

X4 declara quiebra y detiene la simulación si el patrimonio cae a cero o menos, o si hay
margen usado y el nivel de margen cae por debajo de \code{STOP_OUT_LEVEL} (50\,\% por
defecto). Se registra un evento \code{stop_out} con la foto completa de la cuenta, se marca
\code{estado['stop_out'] = True} en el checkpoint y se corta el bucle.

\begin{nota}
En modo X5 el capital inicial es de un millón de dólares virtuales por activo,
deliberadamente desproporcionado. La razón: en recolección de datos no interesa saber si la
cuenta sobrevive, sino observar cómo se comporta cada combinación de parámetros. Un tope de
capital cortaría las pasadas antes de tiempo y sesgaría el conjunto de datos hacia
configuraciones conservadoras.
\end{nota}

\subsubsection{Los dos archivos de patrimonio}

\code{equity_global.csv} recibe una fila por vela con saldo, patrimonio, margen usado,
margen libre, nivel de margen y número de posiciones abiertas y en espera.
\code{equity_activos.csv} recibe una fila por vela \textit{y activo} con tres columnas:
ganancia cerrada acumulada (GC), ganancia flotante actual (GA) y su suma (GT). El primero
sirve para la curva de capital; el segundo para atribuir el resultado a cada activo.

\subsection{Recálculo de soportes dentro del backtest}

\subsubsection{Delegación en X0}

X4 no reimplementa el algoritmo de soportes: importa \code{_procesar_valor_N} de X0 y lo
llama con \code{fecha_hora_max = ts}. Esa única bandera cambia tres comportamientos dentro
de X0: recorta los datos, lee y escribe en el caché de backtesting en vez de en el archivo
de producción, y omite la generación de gráficos.

\subsubsection{Paralelismo por activo}

Los recálculos de distintos activos son independientes, así que se reparten entre hasta
cuatro procesos con un \code{ProcessPoolExecutor}\footnote{\textbf{Proceso} vs.\
\textbf{hilo}: Python no permite que varios hilos ejecuten código Python a la vez (por el
bloqueo global del intérprete o GIL), así que para trabajo intensivo en cálculo hay que usar
procesos separados, cada uno con su propio intérprete. \code{ProcessPoolExecutor} reparte
tareas entre procesos; \code{ThreadPoolExecutor}, entre hilos —útil solo cuando el trabajo
es esperar a que termine otra cosa, como en la orquestación de X5.}. La búsqueda de soportes
es puro cálculo numérico, de modo que aquí los procesos son la elección correcta.

\begin{pseudo}{\code{_recalcular_soportes(estado, ts, cfg, x5_mode)}}
\State \textbf{params_soporte} $\gets \{K, N_{EXP}, \Lambda\}$ del tramo actual si \code{x5_mode}, si no \textit{nada}
\State \textbf{tareas} $\gets$ una por activo, con: activo, $N$, rutas, \textbf{ts}, precios de sus posiciones abiertas
\State \textbf{infos} $\gets$ ejecutar \textbf{tareas} en paralelo (hasta 4 procesos) llamando a \code{X0._procesar_valor_N}
\State informar por consola el rango de precios usado, el arranque en caliente y la FO de cada búsqueda
\For{cada \textbf{activo}}
  \State \textbf{nuevos} $\gets$ leer del caché de backtesting la solución más reciente con $t^* \le$ \textbf{ts}
  \If{\textbf{nuevos} no está vacío} \State soportes del activo $\gets$ ordenar(\textbf{nuevos}) \EndIf
\EndFor
\State \code{ts_ultimo_recalculo} $\gets$ \textbf{ts}
\end{pseudo}

\subsubsection{El caché de backtesting y el arranque en caliente}

Optimizar $N$ soportes desde cero tarda minutos. Como cada recálculo es una variación
pequeña del anterior —cinco días más de datos—, empezar desde la solución previa ahorra la
mayor parte del trabajo. Eso es un \textit{arranque en caliente}\footnote{\textbf{Arranque
en caliente} (\textit{warm start}): iniciar un algoritmo iterativo desde una solución
conocida y razonablemente buena en lugar de un punto aleatorio. Reduce el número de
iteraciones necesarias, a costa de sesgar la búsqueda hacia el entorno de esa solución.}.

\bigsep
El caché \code{\{ACTIVO\}_\{N\}_bt.json} es un diccionario de fecha a lista de soportes. Al
resolver el instante $t$, X0 busca la clave más reciente con $t^* \le t$ y usa esa solución
como punto de partida. Al terminar, añade la solución nueva bajo la clave $t$.

\bigsep
Hay una distinción que suena parecida pero es distinta, y conviene fijarla:

\begin{itemize}[leftmargin=1.4em,itemsep=2pt]
\item \code{warm_start} controla \textbf{de dónde sale la solución inicial}: del caché, o de
puntos aleatorios.
\item \code{cold_start} controla \textbf{si se hereda el paso adaptativo} \code{delta_inicial}
del optimizador. En modo X5 se activa (es decir, no se hereda) porque cada tramo cambia $K$,
$N_{EXP}$ y $\Lambda$: heredar la presión acumulada sobre una función objetivo que ya no es
la misma dejaría al optimizador convergido de entrada, sin mover un solo soporte.
\end{itemize}

\subsubsection{Las posiciones abiertas son anclas}

A la búsqueda de soportes se le pasa la lista de precios de las posiciones abiertas del
activo (\code{ordenes_abiertas_bt}). X0 las trata como fijas: puede mover los demás soportes,
pero no esos. Tiene sentido operativo —hay dinero comprometido en ese nivel exacto, moverlo
en la lista no cambia la realidad— y evita que el paso A cancele órdenes que en realidad ya
son posiciones.

\subsection{Persistencia de operaciones y eventos}

\subsubsection{Los dos registros}

\code{trades.json} contiene una entrada por operación cerrada, con todo lo necesario para
analizarla después: identificador, activo, $N$, nivel del soporte, fechas de apertura y
cierre, precios, motivo del cierre, retorno en dólares y en porcentaje del capital,
duración en velas, peor pérdida flotante y mejor ganancia flotante alcanzadas, capital de la
cuenta al abrir, si se resolvió intra-vela, y una copia del bloque de parámetros vigentes.

\bigsep
\code{events.json} es más granular: registra también las creaciones y cancelaciones de
órdenes y cada movimiento de stop. Sirve para reconstruir la película completa cuando algo
no cuadra.

\subsubsection{Identificadores reconstruibles}

El identificador de una operación se compone como
\code{\{activo\}_\{fecha_apertura\}_\{precio_apertura\}}. No se almacena en ningún lado: se
reconstruye cuando hace falta, tanto al abrir la posición como al cerrarla, y sale idéntico
las dos veces porque depende solo de datos que no cambian. Es una forma barata de tener
identificadores estables sin llevar un contador global ni guardar estado extra.

\subsubsection{Escritura segura}

Las operaciones y eventos se acumulan en listas en memoria y se vuelcan periódicamente. El
volcado lee el archivo existente, le añade lo nuevo, escribe a un archivo temporal y lo
renombra sobre el destino. El renombrado es atómico\footnote{\textbf{Escritura atómica}:
técnica para que un archivo nunca quede a medio escribir. Se escribe primero una copia
temporal completa y luego se renombra sobre el original; el sistema operativo garantiza que
el renombrado ocurre entero o no ocurre. Si el programa muere a mitad, el archivo original
sigue intacto.}: si el proceso muere a mitad, el archivo anterior queda intacto.

\bigsep
La lectura, además, reintenta hasta diez veces con pausas de 0,2 segundos. No es paranoia
gratuita: el proyecto vive en una carpeta sincronizada por OneDrive en Windows, y el
servicio de sincronización bloquea temporalmente los archivos recién escritos. Sin el
reintento, el volcado siguiente fallaba con un error de permisos.

\subsection{Modo X5: X4 al servicio de la recolección}
\label{sec:x4-x5mode}

Todo lo anterior describe X4 como simulador. Esta subsección describe qué se le añade
cuando corre con \code{--x5}. Son cuatro cosas.

\subsubsection{Uno: los parámetros cambian durante la corrida}
\label{sec:x4-x5params}

En modo versión, $K$, $N_{EXP}$, $\Lambda$, $A$, $B$, $N$ y \code{PERDIDA_MAX} son
constantes. En modo X5 se regeneran cada \code{delta_recalculo_soportes} días simulados,
acoplados al recálculo de soportes: primero se eligen los parámetros nuevos, después se
recalculan los soportes con ellos.

\begin{pseudo}{\code{_seleccionar_params_x5} — explorar o explotar}
\For{cada \textbf{activo}}
  \State (\textbf{tipo}, \textbf{modelo}) $\gets$ modelo cargado al inicio del ciclo
  \State \textbf{es_explore} $\gets$ (\textbf{modelo} no existe) \textbf{o} (aleatorio$[0,1) <$ \code{EXPLORATION_RATE})
  \If{\textbf{es_explore}}
     \State parámetros $\gets$ sorteo uniforme dentro de \code{X5_PARAM_RANGES} del activo
     \State $B \gets \alpha \cdot A$ con $\alpha \sim U(0{,}3;\,0{,}7)$ \Comment{$B$ debe quedar por debajo de $A$}
  \Else
     \State \textbf{ctx} $\gets$ contexto de mercado al instante \textbf{ts} (features X3, X2, temporales, cartera)
     \State parámetros $\gets$ \code{X5.inferir_con_contexto}(activo, tipo, modelo, \textbf{ctx}, rangos)
  \EndIf
\EndFor
\State \Return parámetros y etiqueta (\code{EXPLORE} / \code{EXPLOIT}) por activo
\end{pseudo}

\bigsep
Esta es la disyuntiva clásica entre \textbf{explorar} (probar combinaciones nuevas para
aprender) y \textbf{explotar} (usar lo que el modelo ya cree que es mejor). La proporción la
fija \code{EXPLORATION_RATE} = 0,30: incluso con modelo entrenado, tres de cada diez tramos
se sortean al azar.

\begin{nota}
Mantener exploración permanente no es un capricho. Si X4 solo usara las recomendaciones del
modelo, el conjunto de datos se llenaría de filas parecidas entre sí y el modelo aprendería
cada vez mejor una región cada vez más chica del espacio, sin manera de descubrir que otra
región es mejor. Es el sesgo de selección observacional: solo se observa lo que se decidió
observar.
\end{nota}

\bigsep
El muestreo de $B$ merece un comentario. $A$ es la ganancia mínima que activa el stop y $B$ la
distancia a la que el stop sigue al precio. Si $B > A$, el primer stop quedaría por debajo
del precio de entrada y podría cerrar la posición en pérdida pese a haber alcanzado la
ganancia de activación. Sortearlos de forma independiente produciría muchas combinaciones
incoherentes; por eso $B$ se sortea como una fracción de $A$ entre 0,3 y 0,7.

\bigsep
Una vez elegidos, \code{_aplicar_params_x5} los inyecta mutando el objeto de configuración
en memoria. Los escalares ($K$, $N_{EXP}$, $\Lambda$, $A$, $B$, pérdida máxima) toman el
valor del primer activo de la lista; los que son por activo ($N$ y lote) se asignan uno a
uno. Como en modo X5 cada proceso corre un solo activo, ese colapso a escalar es exacto, no
una aproximación.

\subsubsection{Dos: cada operación cerrada se convierte en una fila de entrenamiento}

En cada vela y activo se calcula un fotograma de contexto (\code{_compute_x5_snapshot}):
indicadores técnicos de X3 sobre los datos hasta ese instante, y score fundamental de X2
correspondiente a esa fecha. Ese fotograma se adjunta a la orden en dos momentos:

\begin{itemize}[leftmargin=1.4em,itemsep=2pt]
\item cuando se \textbf{crea} la orden en espera, como bloque \code{x3_oe} / \code{x2_oe};
\item cuando se \textbf{ejecuta} y pasa a posición abierta, como bloque \code{x3_oa} /
\code{x2_oa}.
\end{itemize}

\bigsep
Al cerrarse la posición, \code{_construir_fila_oc} arma la fila completa: parámetros
vigentes, features temporales de los dos momentos, los dos bloques de contexto, el estado de
la cartera al cerrar y, finalmente, los objetivos que el modelo debe aprender a predecir.

\begin{pseudo}{\code{_construir_fila_oc} — de operación cerrada a fila de entrenamiento}
\Require posición cerrada, registro de operación, estado del activo, configuración
\State \textbf{fila} $\gets \{$\code{tipo_registro}: \code{'oc'}, activo, fechas de OE, OA y OC$\}$
\State añadir parámetros vigentes: $N$, $K$, $N_{EXP}$, $\Lambda$, $A$, $B$, multiplicador de lote, pérdida máxima
\State añadir features temporales del instante de la OE (hora, día, mes, festivos, indicadores binarios)
\State añadir \code{x2_*} y \code{x3_*} del instante de la OE
\State añadir contexto de cartera: número de posiciones abiertas y en espera, exposición, media y desviación de retornos abiertos, media de las últimas 5 operaciones cerradas, resultado flotante
\State añadir features temporales, \code{x2_*} y \code{x3_*} del instante de la OA, con sufijo \code{_oa}
\State \textbf{objetivos}: \code{retorno_pct}, \code{pnl_flotante_activo}, \code{pnl_cerrado_activo}
\State \Return \textbf{fila}
\end{pseudo}

\subsubsection{Tres: filas periódicas}

Si solo se registraran operaciones cerradas, el conjunto de datos tendría un sesgo grave:
durante una caída prolongada casi nada se cierra —las posiciones quedan abiertas
acumulando pérdida flotante— y el modelo no vería esos períodos. Para cubrirlos, cada
\code{X5_FREQ_REGISTRO_PERIODICO} velas (4 por defecto) con al menos una posición abierta se
escribe una fila \code{'periodico'} con el mismo contexto pero sin datos de OE ni de OA, cuyo
objetivo relevante es el resultado flotante del activo.

\subsubsection{Cuatro: aislamiento por activo}

Con \code{--activo BTCUSD}, X4 reduce la lista de activos a uno y redirige sus carpetas de
trabajo a \code{bt_BTCUSD}. Así, seis procesos pueden correr a la vez sin pisarse el
checkpoint, la curva de patrimonio ni el caché de soportes. Los parámetros que la
configuración de X5 define por activo ($A$, $B$, pérdida máxima) se colapsan al valor de ese
activo, porque el núcleo de X4 los espera como escalares.

\bigsep
La decisión de reiniciar o continuar también cambia en este modo: X4 la toma solo, sin que el
llamador la imponga. Empieza una pasada nueva si no hay checkpoint, si la pasada anterior
llegó al final de los datos, o si terminó en quiebra; retoma el checkpoint si quedó a mitad
de camino, es decir, si el proceso se interrumpió. Así una pasada completa sigue siendo una
unidad independiente —con su propia secuencia de parámetros— sin tirar a la basura horas de
cómputo por una interrupción.

\subsection{Parámetros de X4 y su efecto}

\begin{table}[h]
\centering
\small
\begin{tabularx}{\linewidth}{@{}l l X@{}}
\toprule
\textbf{Parámetro} & \textbf{Valor típico} & \textbf{Qué controla y qué pasa al subirlo} \\
\midrule
\code{capital_inicial} & 3\,000 / 1\,000\,000 & Saldo de partida. En modo X5 se pone absurdamente alto para que el margen nunca corte la pasada. \\
\code{PERDIDA_MAX_BT} & 120 USD & Corte de pérdidas por posición. Subirlo deja correr las perdedoras: menos cierres en rojo pero peores cuando llegan. \\
\code{MARGEN_LIBRE_MIN_BT} & 50 USD & Colchón mínimo de margen. Subirlo hace la simulación más conservadora: rechaza más órdenes. \\
\code{STOP_OUT_LEVEL} & 50\,\% & Nivel de margen al que se declara la quiebra. \\
\code{delta_recalculo_soportes} & 1 (versión) / 5 (X5) & Días entre recálculos. Bajarlo da soportes más frescos y muchísimo más cómputo; en modo X5 también acelera el cambio de parámetros. \\
$A$ & 2--20 USD & Ganancia que activa el stop y distancia mínima para colocar una orden. Subirlo: menos órdenes, más lejos, operaciones más ambiciosas. \\
$B$ & $0{,}3A$--$0{,}7A$ & Holgura del stop. Subirlo tolera más retrocesos antes de cerrar; bajarlo asegura antes y corta más operaciones prematuramente. \\
$N$ & 40--200 & Soportes activos. Subirlo cubre más rango con entradas más finas y fragmenta más el capital. \\
\code{LOTAJES_M} & 1 & Multiplicador del lote mínimo. Fijo en 1 en todo el pipeline de X5. \\
\bottomrule
\end{tabularx}
\caption{Parámetros de X4. Los siete últimos son precisamente los que X5 explora y optimiza.}
\end{table}

\newpage
% ═══════════════════════════════════════════════════════════════════════════
\section{X5 — Cerebro macro}
\label{sec:x5}
% ═══════════════════════════════════════════════════════════════════════════

\subsection{Propósito}

\subsubsection{El problema: una configuración fija para un mercado que no lo es}

Todos los parámetros del sistema —cuántos soportes mantener, cuánta ganancia exigir antes de
proteger, cuánto dejar respirar al stop, cuánta pérdida tolerar— están hoy escritos en
\code{config.py} como números fijos. Funcionan razonablemente en promedio, y ese promedio es
el problema: el mercado de septiembre no es el de marzo, y BTCUSD no se comporta como GOOGL.

\bigsep
La pregunta que da origen a X5 es directa: \textit{dado cómo está el mercado ahora mismo,
¿qué combinación de parámetros habría rendido mejor en situaciones parecidas?} Responderla a
mano es inviable —son ocho parámetros continuos, seis activos y un contexto que cambia todo
el tiempo—, así que se convierte en un problema de aprendizaje.

\subsubsection{La idea: un modelo sustituto}

Lo que de verdad interesa evaluar es una función carísima:
\[
\text{RENDIMIENTO}(\text{contexto},\, \text{parámetros}) \;=\; \text{resultado de operar así en esa situación}
\]
Evaluarla una vez en la realidad cuesta semanas; evaluarla con X4 cuesta minutos u horas.
Optimizarla directamente —probar miles de combinaciones y quedarse con la mejor— es
inviable en vivo.

\bigsep
La salida es entrenar un \textbf{modelo sustituto}\footnote{\textbf{Modelo sustituto}
(\textit{surrogate model}): modelo estadístico barato que imita una función cara de evaluar.
Se ajusta con un conjunto de evaluaciones reales de la función original y luego se usa como
reemplazo para optimizar, porque consultarlo cuesta milisegundos. Es la misma idea detrás
de la optimización bayesiana.}: un modelo que aprende a predecir el resultado sin correr la
simulación. Una vez entrenado, buscar los mejores parámetros es barato: se le pregunta al
modelo miles de veces con distintas combinaciones y se elige la que predice mejor
resultado.

\bigsep
En una frase: \textbf{X4 evalúa la función cara unas cuantas miles de veces; X5 aprende a
imitarla y después la optimiza.}

\subsubsection{El ciclo cerrado}

\begin{table}[h]
\centering
\small
\begin{tabularx}{\linewidth}{@{}r l X@{}}
\toprule
& \textbf{Etapa} & \textbf{Qué ocurre} \\
\midrule
1 & Recolectar & X5 lanza X4 en modo \code{--x5}. X4 simula, cambia parámetros cada pocos días y escribe una fila por operación cerrada en el store del activo. \\
2 & Entrenar & Cuando el store supera el umbral, X5 entrena un modelo por activo con todas sus filas. \\
3 & Inferir & Con el modelo entrenado se busca la combinación de parámetros que maximiza el resultado predicho para el contexto actual. \\
4 & Realimentar & Esa recomendación se usa en los tramos «explotar» de la siguiente pasada de X4 (y, en Fase~2, en el trading real). Vuelve a 1. \\
\bottomrule
\end{tabularx}
\caption{El bucle de X5. Cada vuelta mejora el modelo con datos que el propio modelo ayudó a generar.}
\end{table}

\subsubsection{Las dos fases de despliegue}

\paragraph{Fase 1 (actual, \code{TIPO_EJECUCION = "est"}).} Todo el bucle vive en simulación.
X4 llena los stores, X5 entrena y escribe sus recomendaciones en
\code{config/active_parameters.json}. El sistema en vivo (X1) \textbf{no lee ese archivo}:
sigue con los valores fijos de \code{config.py}. Es la fase de acumulación de evidencia.

\paragraph{Fase 2 (\code{TIPO_EJECUCION = "din"}).} X1 lee
\code{active_parameters.json} y opera con los parámetros recomendados por activo, y cada
operación real cerrada alimenta directamente el store. El bucle se cierra con datos reales.
La transición es \textbf{por activo}: un activo cuyo estado sea \code{untrained} cae de vuelta
a los valores de \code{config.py} de forma automática, así que BTCUSD puede estar en modo
dinámico mientras TSLA sigue en estático.

\subsection{Interfaz de línea de comandos}

X5 tiene tres modos mutuamente excluyentes y una casilla opcional:

\begin{table}[h]
\centering
\small
\begin{tabularx}{\linewidth}{@{}l X@{}}
\toprule
\textbf{Invocación} & \textbf{Efecto} \\
\midrule
\code{--recolectar} & Genera datos con X4 y entrena solo. Pregunta al inicio si recolectar para todos los activos en paralelo o para uno en recorrido guiado. \\
\code{--recolectar --demo} & Recorrido guiado de un activo sobre un store y un modelo aislados con sufijo \code{_demo}, con pausas explicativas. No toca los datos reales. \\
\code{--infer} & Recomienda parámetros con los modelos actuales y escribe \code{active_parameters.json}. \\
\code{--infer --train} & Reentrena desde el store y después recomienda. \\
\code{--status} & Diagnóstico: cuántas operaciones cerradas hay por activo, qué modelo corresponde y qué dice el archivo de parámetros activos. \\
\code{--activo BTCUSD} & Restringe la operación a un activo. \\
\bottomrule
\end{tabularx}
\caption{Modos de X5. \code{--train} no es un modo: es una casilla que solo tiene efecto junto a \code{--infer}.}
\end{table}

\subsection{El store: el conjunto de datos}

\subsubsection{Un archivo por activo}

\code{resources/x5/\{ACTIVO\}_store.csv} acumula todas las filas generadas para ese activo, a
lo largo de todas las pasadas de todas las corridas. No se sobrescribe: crece. Cada fila es
una observación de la forma \textit{«en este contexto, con estos parámetros, pasó esto»}.

\bigsep
La separación por activo es deliberada: cada activo entrena su propio modelo. Un modelo
único para los seis tendría más datos, pero mezclaría regímenes de volatilidad
incomparables (BTCUSD se mueve varios puntos porcentuales en horas; GOOGL no).

\subsubsection{Los dos tipos de fila}

\begin{itemize}[leftmargin=1.4em,itemsep=3pt]
\item \code{tipo_registro = 'oc'} — se escribe cuando una posición se cierra. Es la
observación completa: contexto en el momento de crear la orden, contexto en el momento de
ejecutarla, parámetros vigentes y resultado.
\item \code{tipo_registro = 'periodico'} — se escribe cada cuatro velas mientras haya al
menos una posición abierta. No tiene datos de OE ni de OA; su valor está en registrar el
resultado flotante durante períodos en que nada se cierra.
\end{itemize}

\subsubsection{Los bloques de columnas}

El orden de las columnas está fijado por \code{_feature_cols_de_store} y se guarda junto al
modelo. Esto es esencial: si mañana el store gana una columna nueva, un modelo entrenado
antes seguiría leyendo su vector en el orden con el que fue entrenado.

\begin{table}[h]
\centering
\small
\begin{tabularx}{\linewidth}{@{}r l l X@{}}
\toprule
\textbf{\#} & \textbf{Bloque} & \textbf{Prefijo} & \textbf{Contenido} \\
\midrule
1 & Parámetros & --- & \code{n_ejecucion}, \code{K}, \code{N_EXP}, \code{LAMBDA}, \code{A}, \code{B}, \code{LOTAJES_M}, \code{PERDIDA_MAX}. Las ocho variables de decisión. \\
2 & Temporal OE & --- & Hora, día de la semana, día del mes, mes; indicadores binarios de día y mes; días hasta y desde el próximo festivo bursátil de EE.UU. \\
3 & Fundamental OE & \code{x2_} & Score global, score comparativo entre activos, score de tendencia y componentes desagregados. \\
4 & Técnico OE & \code{x3_} & Medias móviles, RSI, MACD, ATR, bandas de Bollinger, momento, volatilidad, caída desde máximos, pendiente de tendencia, distancia y densidad de soportes. \\
5 & Cartera & --- & Posiciones abiertas y en espera, exposición en dólares, media y desviación de los retornos abiertos, media de las últimas cinco operaciones cerradas, resultado flotante. \\
6--8 & Contexto OA & sufijo \code{_oa} & Los bloques 2, 3 y 4 repetidos, tomados en el instante en que la orden se ejecutó. \\
\bottomrule
\end{tabularx}
\caption{Los bloques del vector de entrada, en orden. Suman del orden de 240 columnas.}
\end{table}

\bigsep
Que el contexto aparezca dos veces —al crear la orden y al ejecutarla— tiene una razón
concreta: entre esos dos momentos pueden pasar horas o semanas, y lo que ocurrió en medio
importa. Una orden colocada en calma que se ejecuta durante un desplome es una operación
muy distinta de una colocada y ejecutada el mismo día.

\subsubsection{Los tres objetivos}

X5 no predice un solo número, sino tres:

\begin{table}[h]
\centering
\small
\begin{tabularx}{\linewidth}{@{}l l X@{}}
\toprule
\textbf{Nombre} & \textbf{Columna} & \textbf{Qué mide} \\
\midrule
retorno & \code{retorno_pct} & Resultado de \textit{esta} operación como fracción del capital. Es el objetivo principal: el que se maximiza en inferencia. \\
flotante & \code{pnl_flotante_activo} & Resultado no realizado del conjunto de posiciones abiertas del activo. Único objetivo disponible en las filas periódicas. \\
cerrado & \code{pnl_cerrado_activo} & Resultado realizado acumulado del activo hasta ese momento. \\
\bottomrule
\end{tabularx}
\caption{Objetivos del modelo. Los tres se entrenan; solo el primero guía la optimización.}
\end{table}

\bigsep
¿Por qué entrenar tres si solo se optimiza uno? Porque los otros dos actúan como
\textbf{tareas auxiliares}: obligan a la parte compartida del modelo a construir
representaciones que expliquen también la salud global de la cartera, no solo el desenlace
de una operación aislada. En el modelo neuronal esto es explícito (un tronco compartido y
tres cabezas de salida) y su peso en la función de pérdida es menor: $l_{total} = l_{ret} +
0{,}3\, l_{flot} + 0{,}3\, l_{cerr}$.

\subsubsection{Escritura desde X4}

\code{_append_x5_store} añade la fila al CSV en modo \textit{append}, escribiendo la
cabecera solo si el archivo no existía. Antes de escribir comprueba que el archivo termine
en salto de línea: si una ejecución anterior se cortó a mitad de fila, sin ese control la
fila nueva se concatenaría a la truncada y corrompería el archivo. En la lectura, el
complemento de esa defensa es \code{on_bad_lines='skip'}: una fila malformada se descarta en
vez de abortar la carga completa del store.

\subsection{Recolección}

\subsubsection{El orquestador}

\begin{pseudo}{\code{_recolectar()} — recolección paralela}
\State \textbf{cfg_x5} $\gets$ cargar \code{scripts/config_x5_default.py}
\State \textbf{activos} $\gets$ \code{cfg_x5.valores}; \quad \textbf{n_ciclos} $\gets$ \code{cfg_x5.N_CICLOS_BT}
\If{hay terminal interactiva}
  \For{cada \textbf{activo}} preguntar si reiniciarlo desde cero; si sí, borrar todo lo suyo \EndFor
\EndIf
\State imprimir estado inicial: operaciones cerradas por activo
\State lanzar un hilo por activo con \code{ThreadPoolExecutor}
\For{cada worker que termina}
  \State informar cuántas operaciones cerradas ganó ese activo
\EndFor
\State mostrar el estado final (\code{--status})
\end{pseudo}

\bigsep
El paralelismo aquí es de \textbf{hilos}, no de procesos, y no contradice lo dicho en la
Sección~\ref{sec:x4-x5mode}: cada hilo no calcula nada, solo lanza un subproceso de X4 y
espera a que termine. El trabajo pesado ocurre en esos seis procesos hijos independientes.

\subsubsection{El worker de un activo}

\begin{pseudo}{\code{_worker_recolectar_activo(activo, x4_path, n_ciclos)}}
\State \textbf{n_ini} $\gets$ operaciones cerradas en el store del activo
\For{\textbf{ciclo} $= 1 \dots$ \textbf{n_ciclos}}
  \State ejecutar \code{python X4_backtester.py --x5 --activo \{activo\}} capturando su salida
  \State \textbf{n} $\gets$ operaciones cerradas en el store
  \If{el subproceso falló} imprimir las últimas líneas de error y \textbf{cortar} \EndIf
  \If{\textbf{n} $\ge$ \code{X5_MIN_TRADES_TRAIN}}
     \State \textbf{entrenar}(activo) \Comment{los tramos de explotación del ciclo siguiente ya usan este modelo}
  \EndIf
  \If{\textbf{n} $\ge$ \code{X5_MIN_TRADES_FTT}} \textbf{cortar} \Comment{ya hay datos de sobra} \EndIf
\EndFor
\State \Return activo y su recuento final
\end{pseudo}

\bigsep
Cada iteración de este bucle es un \textbf{ciclo}. Un ciclo lanza X4 una vez; X4 decide por
su cuenta si eso significa empezar una pasada nueva desde \code{fecha_inicio} o retomar una
interrumpida. La consecuencia práctica: el reentrenamiento ocurre \textbf{entre} ciclos,
nunca durante uno. Dentro de un ciclo el modelo está congelado —se carga una sola vez al
principio, en \code{ejecutar_x5_ciclo}—, lo cual mantiene la simulación coherente: todos los
tramos de explotación de una pasada consultan la misma versión del modelo.

\subsubsection{El recorrido guiado}

\code{_recolectar_guiado} corre los ciclos de forma \textbf{secuencial} —no paralela— con la
salida de X4 heredada, de modo que un narrador (\code{x5_demo.py}) pueda explicar cada
suceso nuevo la primera vez que aparece. Los sucesos tienen identificadores estables:

\begin{table}[h]
\centering
\small
\begin{tabularx}{\linewidth}{@{}l l X@{}}
\toprule
\textbf{ID} & \textbf{Tipo} & \textbf{Suceso} \\
\midrule
D01 & recurrente & Recálculo de soportes (arranque en frío o periódico) \\
D02 & recurrente & Tramo de exploración: parámetros sorteados al azar \\
D03 & recurrente & Orden en espera creada \\
D04 & recurrente & Orden ejecutada: pasa a posición abierta \\
D05 & recurrente & Posición cerrada: se añadió una fila al store \\
D06 & hito & Tramo de explotación: primera inferencia con el modelo \\
D07 & recurrente & Avance de mes dentro del backtest \\
D08 & hito & Se cruzó el umbral de entrenamiento; se entrena el modelo \\
D09 & hito & Fin de la recolección del activo \\
D10 & recurrente & Gráfico de la búsqueda de soportes recién guardado \\
\bottomrule
\end{tabularx}
\caption{Sucesos del narrador. Los marcados como hito vuelven a pausar aunque ya se hayan visto.}
\end{table}

\bigsep
El mismo mecanismo tiene dos usos según dos variables de entorno:

\begin{itemize}[leftmargin=1.4em,itemsep=2pt]
\item \code{--recolectar --demo}: sufijo \code{_demo} activo. El store, el modelo, el
historial de métricas, el checkpoint del backtest y los gráficos van a rutas paralelas. Se
pausa en cada suceso nuevo. Los datos reales quedan intactos.
\item \code{--recolectar} con alcance «uno»: sin sufijo. Misma narración y mismos gráficos,
pero sobre los datos reales y \textbf{sin pausas} —una recolección oficial no puede quedarse
esperando que alguien pulse una tecla.
\end{itemize}

\subsubsection{Reiniciar un activo}

\code{_borrar_checkpoints_activo} borra absolutamente todo lo asociado al activo dentro de su
espacio de nombres: store, estado del narrador, carpeta del modelo, historial de métricas,
carpeta completa del backtest (checkpoint, curvas de patrimonio, caché de soportes, registros)
y gráficos. El significado es literal: como si el activo nunca se hubiera procesado.

\begin{nota}
El mismo alcance aplica en el espacio de nombres real, donde el borrado destruye datos
acumulados durante horas de cómputo. La pregunta de reinicio se hace explícitamente al
inicio y por activo, y solo si hay terminal interactiva: en ejecución automatizada no se
pregunta nada y no se borra nada.
\end{nota}

\bigsep
Cada borrado se hace con reintentos y espera creciente (0,5\,s, 1\,s, 2\,s, 4\,s, 5\,s\dots
hasta unos 35 segundos en total). En Windows, OneDrive y el antivirus bloquean archivos
recién escritos durante bastante más que un par de segundos, y sin ese reintento el
reinicio fallaba con un error de permisos.

\subsection{Entrenamiento}

\subsubsection{Qué modelo, según cuántos datos haya}

X5 elige el modelo por volumen de datos, contando solo filas de tipo \code{'oc'}:

\begin{table}[h]
\centering
\small
\begin{tabularx}{\linewidth}{@{}l l X@{}}
\toprule
\textbf{Operaciones cerradas} & \textbf{Modelo} & \textbf{Razón} \\
\midrule
$< 500$ & \code{untrained} & Sin datos suficientes. Se devuelven los valores de \code{config.py} tal cual. \\
$500 \dots 4999$ & LightGBM & Los métodos basados en árboles son muy eficientes con pocos datos tabulares y no necesitan normalización. \\
$\ge 5000$ & FT-Transformer & Con suficientes datos, una red neuronal captura interacciones que los árboles aproximan a saltos. \\
\bottomrule
\end{tabularx}
\caption{Selección automática de modelo (\code{_seleccionar_tipo}).}
\end{table}

\subsubsection{Preparación del vector de entrada}

\begin{pseudo}{\code{_preparar_features(store, columnas, objetivo, tipos)}}
\Require store completo, lista ordenada de columnas, nombre del objetivo, tipos de fila admitidos
\Ensure matriz $X$, vector $y$, vector de pesos $w$
\State \textbf{df} $\gets$ filas del store cuyo \code{tipo_registro} esté en \textbf{tipos}
\State descartar filas cuyo objetivo no sea numérico \Comment{p.ej.\ \code{retorno_pct} vacío en filas periódicas}
\State $X \gets$ matriz $|df| \times |$columnas$|$; cada columna convertida a número, ausentes y no numéricos a $0$
\State $y \gets$ columna objetivo como número
\If{\code{X5_WINDOW_TRAIN} está definido} quedarse con las últimas \code{X5_WINDOW_TRAIN} filas \EndIf
\State $w_i \gets \exp\!\big(\lambda_{decay}\,(i - n + 1)\big)$ para $i = 0 \dots n-1$
\State $w \gets w / \overline{w}$ \Comment{media 1, para no alterar la escala de la pérdida}
\State \Return $X,\, y,\, w$
\end{pseudo}

\paragraph{La ponderación temporal.} La última fila del store pesa $\approx 1$ y las
anteriores decaen exponencialmente hacia atrás con $\lambda = 0{,}001$; a unas 700 filas de
distancia el peso es la mitad. La intención es que el modelo escuche más al pasado reciente
sin descartar el antiguo. La normalización posterior a media 1 evita que el tamaño del store
cambie la magnitud efectiva de la función de pérdida.

\paragraph{Los ausentes van a cero.} Es una decisión discutible y conviene tenerla presente:
una columna vacía y una columna que vale cero son indistinguibles para el modelo. Con los
árboles el efecto es menor; en el modelo neuronal, un cero tras normalizar cae en la media de
la columna, que es un valor neutro razonable.

\subsubsection{LightGBM: tres modelos independientes}

LightGBM\footnote{\textbf{LightGBM}: implementación de \textit{gradient boosting} sobre
árboles de decisión. El método construye cientos de árboles pequeños en secuencia, donde cada
árbol nuevo se entrena para corregir el error que dejaron los anteriores. Es el estándar de
facto en datos tabulares por su precisión y velocidad, y no requiere normalizar las
variables.} entrena un modelo separado por objetivo, guardados como tres archivos distintos
más un cuarto con la lista ordenada de columnas.

\begin{pseudo}{\code{_entrenar_lgbm(activo, store)}}
\State \textbf{columnas} $\gets$ orden canónico de features del store
\For{cada \textbf{objetivo} en \{retorno (solo \code{oc}), flotante (\code{oc} y \code{periodico}), cerrado (solo \code{oc})\}}
  \State $X, y, w \gets$ \textbf{preparar_features}(store, \textbf{columnas}, \textbf{objetivo}, tipos)
  \If{$|X| < 20$} avisar y \Omitir \EndIf
  \State \textbf{corte} $\gets \max(10,\; \lfloor 0{,}8\,|X| \rfloor)$ \Comment{división temporal, sin mezclar}
  \State entrenar con $X[:corte]$ ponderado por $w$, validando en $X[corte:]$
  \If{el conjunto de validación tiene al menos 10 filas} activar parada temprana con paciencia 50 \EndIf
  \State calcular métricas en entrenamiento y validación; guardar el modelo
\EndFor
\State guardar la lista de columnas y el historial de métricas
\end{pseudo}

\paragraph{La división es temporal, no aleatoria.} El primer 80\,\% de las filas —las más
antiguas— entrena; el último 20\,\% valida. Barajar sería un error: dejaría filas del futuro
en entrenamiento y filas del pasado en validación, y la métrica resultante sería
optimista sin serlo de verdad.

\paragraph{Parada temprana.} Se entrenan hasta mil árboles, pero si el error de validación no
mejora durante cincuenta rondas seguidas, se detiene y se conserva la mejor
iteración\footnote{\textbf{Parada temprana} (\textit{early stopping}): mecanismo para evitar
el sobreajuste. Se vigila el error sobre datos que el modelo no está usando para aprender; en
cuanto ese error deja de mejorar, se detiene el entrenamiento aunque quedaran rondas por
hacer. \textbf{Sobreajuste} (\textit{overfitting}): cuando el modelo memoriza el ruido del
conjunto de entrenamiento y por eso rinde peor con datos nuevos.}.

\paragraph{El objetivo flotante usa más filas.} Es el único que se entrena con filas
\code{'oc'} y \code{'periodico'} juntas, porque es el único objetivo presente en ambos tipos.
Ahí es donde entra la información de los períodos bajistas prolongados en que nada se cierra.

\subsubsection{FT-Transformer: un tronco y tres cabezas}

Cuando el store supera las 5\,000 operaciones cerradas se pasa a una red neuronal de tipo
\textit{Feature Tokenizer + Transformer}, con los hiperparámetros del artículo original
(Gorishniy et al., 2021): dimensión interna $d = 192$, tres bloques, ocho cabezas de atención
y red interna de 256 unidades. Con unas 240 columnas de entrada resultan alrededor de 700\,000
parámetros, que se entrenan sin dificultad en una máquina normal.

\bigsep
La arquitectura tiene tres partes:

\paragraph{Tokenizador de features.} Cada columna es un solo número, y un
\textit{transformer}\footnote{\textbf{Transformer}: arquitectura de red neuronal basada en el
mecanismo de \textbf{atención}, que permite a cada elemento de una secuencia mirar a todos los
demás y decidir a cuáles hacer caso. Se popularizó en procesamiento de lenguaje —cada palabra
atiende al resto de la frase— y aquí se aplica a columnas de una tabla: cada variable puede
condicionar su efecto al valor de las otras.} trabaja sobre secuencias de vectores. El
tokenizador convierte cada número escalar $x_j$ en un vector de dimensión $d$ mediante pesos
propios de esa columna: $t_j = x_j \cdot W_j + b_j$. Así, «RSI vale 30» y «hora vale 30» se
convierten en vectores completamente distintos.

\paragraph{Bloques transformer.} Tres bloques con normalización previa: normalización de
capa\footnote{\textbf{Normalización de capa} (\textit{LayerNorm}): reescala los valores dentro
de cada vector para que tengan media cero y desviación uno. Estabiliza el entrenamiento
evitando que las magnitudes crezcan o se desvanezcan al atravesar muchas capas. La variante
\textit{pre-norm} —normalizar antes de cada subcapa, no después— es la que mejor se comporta
en redes profundas.}, atención de ocho cabezas, conexión residual; luego normalización, red
de dos capas con activación GELU\footnote{\textbf{GELU}: función de activación suave, variante
moderna de ReLU. Introduce la no linealidad sin la que la red colapsaría a una simple
combinación lineal.}, conexión residual.

\paragraph{Agregación y cabezas.} Los 240 vectores de salida se promedian en uno solo
(\textit{mean pooling}) y de ahí salen tres capas lineales, una por objetivo.

\begin{pseudo}{\code{_entrenar_ftt(activo, store)}}
\State $X, y_{ret}, w \gets$ \textbf{preparar_features}(store, columnas, \code{retorno_pct}, (\code{oc}))
\If{$|X| < 50$} avisar y \Return falso \EndIf
\State $y_{flot},\, y_{cerr} \gets$ los otros dos objetivos sobre las mismas filas
\State \textbf{corte} $\gets \max(10, \lfloor 0{,}8 n \rfloor)$
\State \textbf{media}, \textbf{desv} $\gets$ estadísticos por columna calculados \textbf{solo} sobre $X[:corte]$
\State $X_{norm} \gets (X - \textbf{media}) / \textbf{desv}$ \Comment{se guardan para poder normalizar igual en inferencia}
\State inicializar la red y el optimizador AdamW (tasa $10^{-3}$, regularización $10^{-4}$)
\For{\textbf{época} $= 1 \dots 50$}
  \For{cada lote de hasta 256 filas, en orden aleatorio}
     \State $(\hat y_{ret}, \hat y_{flot}, \hat y_{cerr}) \gets$ red$(X_{lote})$
     \State $l \gets \overline{w(\hat y_{ret} - y_{ret})^2} + 0{,}3\,\overline{w(\hat y_{flot} - y_{flot})^2} + 0{,}3\,\overline{w(\hat y_{cerr} - y_{cerr})^2}$
     \State retropropagar, recortar el gradiente a norma 1, dar un paso del optimizador
  \EndFor
\EndFor
\State guardar pesos, número de columnas y estadísticos de normalización; calcular métricas
\end{pseudo}

\bigsep
Dos detalles del pseudocódigo merecen énfasis:

\begin{itemize}[leftmargin=1.4em,itemsep=2pt]
\item \textbf{Los estadísticos de normalización se calculan solo con el tramo de
entrenamiento} y se guardan junto al modelo. Calcularlos sobre todo el conjunto sería una
fuga de información sutil: la media del conjunto de validación entraría en el
preprocesamiento.
\item \textbf{El recorte de gradiente a norma 1} evita que un lote atípico —una operación con
un retorno extremo— produzca un paso gigante que desestabilice el entrenamiento.
\end{itemize}

\begin{nota}
Limitación conocida: como las tres cabezas comparten el mismo tensor de entrada, el
FT-Transformer entrena los tres objetivos únicamente sobre filas \code{'oc'}. A diferencia de
LightGBM, no aprovecha las filas periódicas para el objetivo flotante. Aprovecharlas exigiría
un paso separado, y está anotado como pendiente.
\end{nota}

\subsubsection{Métricas y su historial}

Tras cada entrenamiento se calculan cuatro métricas por objetivo, sobre entrenamiento y sobre
validación: error absoluto medio (MAE), error cuadrático medio (MSE), error porcentual
absoluto medio (MAPE) y coeficiente de determinación
$R^2$\footnote{\textbf{Coeficiente de determinación} ($R^2$): fracción de la variabilidad del
objetivo que el modelo explica. Vale 1 si acierta perfectamente y 0 si no aporta nada sobre
predecir siempre la media. Puede ser \textbf{negativo}, y en este dominio es lo esperable al
principio: significa que el modelo predice peor que la media constante. El MAPE, por su parte,
se calcula solo sobre filas cuyo valor real no sea prácticamente cero, porque dividir por algo
minúsculo lo dispara.}.

\bigsep
Todo se acumula como historial en
\code{resources/x5/Performance/\{ACTIVO\}_performance.json}: una entrada por entrenamiento,
con marca de tiempo, tipo de modelo y tamaños de los dos conjuntos. Permite ver si el modelo
mejora conforme el store crece, que es la única señal real de que el bucle está funcionando.

\subsection{Inferencia: de un modelo entrenado a una recomendación}

Aquí se invierte la pregunta. En entrenamiento se busca la función; en inferencia la función
está fija y se buscan las entradas que la maximizan:
\[
\text{parámetros}^{*} \;=\; \arg\max_{\text{parámetros} \,\in\, \text{rangos}}
\; \widehat{\text{retorno}}(\text{contexto fijo},\, \text{parámetros})
\]
El contexto es un dato, no una variable: describe el mercado tal como está. Las ocho columnas
de parámetros son las únicas que se mueven.

\subsubsection{El contexto}

Hay dos formas de construirlo, con el mismo esquema de columnas:

\begin{itemize}[leftmargin=1.4em,itemsep=2pt]
\item \textbf{En vivo} (\code{_leer_contexto_actual}): features temporales de ahora, último
score fundamental de X2, última fila de features técnicas de X3. Las columnas de cartera se
ponen a cero —en Fase 1 no hay cartera viva que consultar— y las columnas con sufijo
\code{_oa} se rellenan copiando las de OE, bajo el supuesto de que la orden se ejecutará
pronto.
\item \textbf{Dentro del backtest} (\code{_contexto_inferencia_x5}, en X4): idéntico
esquema, pero con datos recortados al instante simulado y con la cartera \textbf{real} de ese
momento.
\end{itemize}

\bigsep
Que ambos produzcan exactamente las mismas claves no es un detalle estético: si difirieran,
el vector de inferencia no encajaría con el orden de columnas con el que se entrenó, y el
modelo leería valores en posiciones equivocadas sin dar ningún error. Por eso las funciones
\code{_extraer_x2} de X4 y de X5 están duplicadas con una advertencia explícita de mantenerlas
sincronizadas.

\subsubsection{Optimización sobre LightGBM: búsqueda bayesiana}

Un modelo de árboles no es derivable, así que no se puede «subir por el gradiente». Se usa
Optuna\footnote{\textbf{Optuna}: biblioteca de optimización de caja negra. Su algoritmo por
defecto, TPE (\textit{Tree-structured Parzen Estimator}), construye a partir de los intentos
ya evaluados dos distribuciones de probabilidad —una de los intentos buenos y otra de los
malos— y propone el siguiente punto donde el cociente entre ambas es más favorable. En la
práctica: concentra los intentos en las zonas prometedoras en vez de repartirlos al azar.},
con 200 intentos por inferencia.

\begin{pseudo}{\code{_inferir_lgbm} — búsqueda bayesiana de parámetros}
\Require modelo del objetivo \code{retorno}, contexto fijo, orden de columnas
\Function{objetivo}{intento}
  \State $p \gets$ sortear los ocho parámetros dentro de sus rangos según lo que sugiera el intento
  \State \Comment{\code{LAMBDA} se sugiere en escala logarítmica: su rango es pequeño y positivo}
  \State \textbf{vector} $\gets$ contexto fijo, con las ocho posiciones de parámetros reemplazadas por $p$
  \State \Return predicción del modelo sobre \textbf{vector}
\EndFunction
\State \textbf{estudio} $\gets$ crear estudio de maximización con muestreador TPE y semilla fija
\State ejecutar \code{X5_N_OPTUNA_TRIALS} intentos
\State \Return los mejores parámetros, redondeados y con los enteros convertidos
\end{pseudo}

\bigsep
La semilla fija (42) hace la recomendación reproducible: mismo contexto y mismo modelo dan
siempre el mismo resultado.

\subsubsection{Optimización sobre el FT-Transformer: ascenso por gradiente}

Una red neuronal sí es derivable respecto de su entrada, y eso permite una técnica más
directa: tratar los parámetros como variables a optimizar y subir por el
gradiente\footnote{\textbf{Ascenso por gradiente}: el gradiente de una función indica la
dirección de máximo crecimiento. Moverse repetidamente en esa dirección lleva a un máximo
local. Es el mismo mecanismo del entrenamiento —que \textit{desciende} por el gradiente para
minimizar el error— pero aplicado a las \textbf{entradas} y con los pesos de la red
congelados.}.

\bigsep
El obstáculo de implementación es que solo ocho de las 240 posiciones del vector deben ser
variables; las otras 232 son constantes. La solución es construir el vector de entrada como
\[
x \;=\; \text{contexto}_{norm} \;+\; P^{\top} p
\]
donde $P$ es una matriz constante de ceros y unos que coloca cada parámetro en su posición, y
$p$ es el vector de ocho variables. Como $P$ es constante, el gradiente fluye limpiamente
hacia $p$ y solo hacia $p$.

\begin{pseudo}{\code{_inferir_ftt} — ascenso por gradiente con reinicios}
\State \textbf{índices} $\gets$ posiciones de las ocho columnas de parámetros dentro del vector
\State \textbf{base} $\gets$ contexto con esas ocho posiciones puestas a cero, normalizado con los estadísticos guardados
\State $P \gets$ matriz de proyección constante; $lo, hi \gets$ rangos de cada parámetro, normalizados
\State \textbf{mejor} $\gets$ ninguno; \quad \textbf{mejor_valor} $\gets -\infty$
\For{\textbf{reinicio} $= 1 \dots$ \code{X5_ASCENT_RESTARTS}}
  \State $p \gets$ punto aleatorio uniforme dentro de $[lo, hi]$
  \For{\textbf{paso} $= 1 \dots$ \code{X5_ASCENT_STEPS}}
     \State $x \gets \textbf{base} + P^{\top} p$
     \State retropropagar $-\hat y_{ret}(x)$ y dar un paso del optimizador Adam
     \State recortar $p$ al rango $[lo, hi]$ \Comment{la búsqueda no puede salirse de lo permitido}
  \EndFor
  \State evaluar $\hat y_{ret}$ en el $p$ final; si supera a \textbf{mejor_valor}, guardarlo
\EndFor
\State desnormalizar \textbf{mejor}, recortar de nuevo a los rangos, redondear y convertir enteros
\State completar con los valores de \code{config.py} cualquier parámetro que no estuviera en el vector
\end{pseudo}

\bigsep
Los diez reinicios desde puntos aleatorios distintos son la defensa contra los máximos
locales: la superficie predicha por una red no es cóncava y un solo ascenso puede quedarse
atrapado en una loma cualquiera. El recorte al rango dentro del bucle —no solo al final—
importa: sin él, el ascenso puede alejarse mucho hacia una zona sin datos de entrenamiento,
donde la predicción del modelo no significa nada, y volver al borde al final daría un
resultado peor que haberse quedado dentro todo el camino.

\subsubsection{El airbag: una regla que el modelo no puede sobrescribir}

Antes de entregar la recomendación en vivo se aplica una regla explícita, no aprendida:

\begin{pseudo}{\code{_aplicar_airbag(activo, params)}}
\State leer los últimos 5 cierres horarios del activo
\State \textbf{caida} $\gets (p_{final} - p_{inicial}) / p_{inicial}$
\If{\textbf{caida} $< -$\code{X5_AIRBAG_THRESHOLD[activo]}} \Comment{8\,\% en cripto, 5\,\% en acciones}
  \State \code{LOTAJES_M} $\gets 1$ \Comment{lote mínimo}
  \State \code{n_sizes_ejecucion} $\gets$ \code{X5_N_MINIMO[activo]} \Comment{30 en cripto, 20 en acciones}
\EndIf
\State \Return parámetros
\end{pseudo}

\bigsep
La justificación es de gestión de riesgo, no de estadística: durante un desplome, los datos
de entrenamiento están fuera de su distribución habitual y la recomendación del modelo no es
confiable; ante la duda, se reduce exposición. Es una capa de seguridad deliberadamente
tonta y deliberadamente inapelable.

\begin{nota}
El airbag \textbf{no} se aplica durante la recolección (\code{inferir_con_contexto} lo omite
explícitamente). El motivo es coherente con lo anterior: en recolección se busca observar qué
pasa realmente en regímenes de caída, incluso con parámetros agresivos. Aplicar el airbag ahí
recortaría precisamente los datos que hacen falta para aprender de esas situaciones.
\end{nota}

\subsubsection{La salida: \code{active_parameters.json}}

\begin{lstlisting}
{
  "generated_at": "2026-08-26T09:12:44",
  "BTCUSD": { "model_status": "lgbm", "n_sizes_ejecucion": 173, "K": 1.2841,
              "N_EXP": 1.9022, "LAMBDA": 0.004112, "A": 8.44, "B": 3.91,
              "LOTAJES_M": 1, "PERDIDA_MAX": 187.4 },
  "TSLA":   { "model_status": "untrained", "n_sizes_ejecucion": 250, ... }
}
\end{lstlisting}

\bigsep
La escritura es un \textbf{mezclado}, no un reemplazo: se lee el archivo existente, se
actualizan solo las claves de los activos recién inferidos y se conservan las demás. Sin eso,
seis procesos paralelos ejecutando \code{--infer --activo X} se borrarían el trabajo unos a
otros. El campo \code{model_status} es el interruptor de seguridad que lee X1: con valor
\code{untrained}, ese activo usa los valores de \code{config.py}.

\subsection{Pseudocódigo consolidado de X5}

\begin{pseudo}{\code{main()} — los tres modos}
\If{\code{--status}} mostrar recuento por activo, modelo correspondiente y contenido del archivo de parámetros; \Return \EndIf
\If{\code{--recolectar}}
  \If{\code{--demo}} preguntar el activo y correr el recorrido guiado sobre el espacio de nombres \code{_demo}
  \ElsIf{el usuario elige «todos»} \textbf{recolectar_paralelo}()
  \Else{} preguntar el activo y correr el recorrido guiado oficial (sin pausas) \EndIf
  \State \Return
\EndIf
\State \Comment{a partir de aquí, \code{--infer}}
\State \textbf{activos} $\gets$ el indicado con \code{--activo}, o todos
\If{\code{--train}} \For{cada activo} \textbf{entrenar}(activo) \EndFor \EndIf
\For{cada \textbf{activo}}
  \State (\textbf{tipo}, \textbf{modelo}) $\gets$ \textbf{cargar_modelo}(activo)
  \If{\textbf{tipo} $=$ \code{untrained}}
     \State parámetros $\gets$ valores de \code{config.py}; estado $\gets$ \code{untrained}
  \Else
     \State \textbf{ctx} $\gets$ contexto de mercado actual
     \State parámetros $\gets$ búsqueda bayesiana o ascenso por gradiente, según \textbf{tipo}
     \State parámetros $\gets$ \textbf{airbag}(activo, parámetros)
  \EndIf
\EndFor
\State escribir \code{config/active_parameters.json} mezclando con lo ya existente
\end{pseudo}

\subsection{Parámetros de X5 y su efecto}

\begin{table}[h]
\centering
\small
\begin{tabularx}{\linewidth}{@{}l l X@{}}
\toprule
\textbf{Parámetro} & \textbf{Valor} & \textbf{Qué controla y qué pasa al subirlo} \\
\midrule
\code{X5_MIN_TRADES_TRAIN} & 500 & Umbral para entrenar. Bajarlo da modelos antes, pero muy ruidosos y con recomendaciones erráticas. \\
\code{X5_MIN_TRADES_FTT} & 5\,000 & Umbral de cambio a red neuronal. Bajarlo arriesga sobreajuste; subirlo posterga un modelo más expresivo. \\
\code{EXPLORATION_RATE} & 0,30 & Fracción de tramos con parámetros al azar. Subirlo da más variedad y peor rendimiento simulado; bajarlo concentra los datos y limita el aprendizaje. \\
\code{N_CICLOS_BT} & 10 & Pasadas de X4 por activo y ejecución de \code{--recolectar}. \\
\code{X5_FREQ_REGISTRO_PERIODICO} & 4 & Velas entre filas periódicas. Bajarlo aumenta mucho el tamaño del store con filas muy correlacionadas entre sí. \\
\code{X5_LAMBDA_DECAY} & 0,001 & Olvido del pasado. Subirlo hace al modelo más reactivo y más olvidadizo. \\
\code{X5_WINDOW_TRAIN} & \code{None} & Tope de filas de entrenamiento. \code{None} = usar todo el store. \\
\code{X5_N_OPTUNA_TRIALS} & 200 & Intentos por inferencia con LightGBM. Subirlo mejora la búsqueda y alarga la inferencia (10--30\,s con 200). \\
\code{X5_ASCENT_RESTARTS} & 10 & Puntos de partida del ascenso por gradiente. Subirlo reduce el riesgo de máximos locales, linealmente más lento. \\
\code{X5_ASCENT_STEPS} & 300 & Pasos por reinicio. \\
\code{X5_ASCENT_LR} & 0,01 & Tamaño del paso. Subirlo explora más rápido y puede oscilar sin converger. \\
\code{X5_AIRBAG_THRESHOLD} & 0,08 / 0,05 & Caída en 4 velas que dispara el airbag. Bajarlo lo hace más sensible. \\
\bottomrule
\end{tabularx}
\caption{Parámetros de X5. Los tres primeros gobiernan el ritmo del bucle completo.}
\end{table}

\newpage
% ═══════════════════════════════════════════════════════════════════════════
\section{Interacción X4 $\leftrightarrow$ X5}
\label{sec:interaccion}
% ═══════════════════════════════════════════════════════════════════════════

\subsection{Quién llama a quién}

La dependencia es mutua y vale la pena verla explícita, porque a primera vista parece
circular sin serlo:

\begin{table}[h]
\centering
\small
\begin{tabularx}{\linewidth}{@{}l l X@{}}
\toprule
\textbf{Desde} & \textbf{Hacia} & \textbf{Mecanismo} \\
\midrule
X5 & X4 & Subproceso: \code{python X4_backtester.py --x5 --activo \{A\}}. Comunicación por archivos y por salida estándar. \\
X4 & X5 & Importación directa: \code{import X5_macro_brain}, para cargar el modelo del activo e inferir parámetros. \\
X4 & X0 & Importación directa de \code{_procesar_valor_N} y \code{_bt_warm_start}. \\
X4 & X3 & Importación de \code{compute_snapshot} para el fotograma de indicadores técnicos. \\
X4 & X2 & Lectura del archivo \code{x2_history.json}, sin ejecutar nada. \\
\bottomrule
\end{tabularx}
\caption{Llamadas entre módulos. Solo la primera cruza una frontera de proceso.}
\end{table}

\bigsep
No hay ciclo de importación porque X4 importa X5 \textit{dentro} de las funciones que lo
necesitan, no en la cabecera del archivo. Y X5 nunca importa X4: lo lanza como programa
externo.

\subsection{Los tres relojes}

Buena parte de la confusión al leer el código viene de que hay tres escalas de tiempo
distintas funcionando a la vez, y las tres se llaman «ciclo» en algún comentario. Conviene
separarlas de una vez:

\begin{table}[h]
\centering
\small
\begin{tabularx}{\linewidth}{@{}l l l X@{}}
\toprule
\textbf{Reloj} & \textbf{Unidad} & \textbf{Duración} & \textbf{Qué cambia en cada tic} \\
\midrule
Vela & 1 hora simulada & --- & Se ejecutan los pasos A--F; se registra el patrimonio. \\
Tramo & \code{delta_recalculo_soportes} días (5) & $\approx$120 velas & Se eligen parámetros nuevos y se recalculan los soportes. \\
Ciclo & Una pasada completa de X4 & Todo el histórico & Se reentrena el modelo del activo (entre ciclos, nunca dentro). \\
\bottomrule
\end{tabularx}
\caption{Las tres escalas temporales del pipeline X4\,+\,X5.}
\end{table}

\bigsep
De estos tres relojes se deduce el orden de magnitud del conjunto de datos. Con datos desde
enero de 2024, un ciclo cubre unas 15\,000 velas horarias por activo, es decir alrededor de
125 tramos, cada uno con su propia combinación de parámetros. Diez ciclos por activo dan
unos 1\,250 experimentos distintos, cada uno generando tantas filas como operaciones cierre.
Ese es el orden de magnitud que hay que alcanzar para cruzar los umbrales de 500 y 5\,000
operaciones cerradas.

\subsection{Una recolección completa, paso a paso}

\begin{pseudo}{Recorrido completo de \code{X5 --recolectar} para un activo}
\State \textbf{[X5]} pregunta el alcance y, si procede, si reiniciar el activo
\State \textbf{[X5]} lanza el worker del activo; empieza el \textbf{ciclo 1}
\State \quad \textbf{[X5$\to$X4]} subproceso \code{X4 --x5 --activo BTCUSD}
\State \quad \textbf{[X4]} carga velas H1 y M1; decide si empieza pasada nueva o retoma checkpoint
\State \quad \textbf{[X4$\to$X5]} \code{cargar_modelo_para_activo(BTCUSD)} $\to$ \code{untrained} en el primer ciclo
\State \quad \textbf{[X4]} sin modelo, \textbf{todos} los tramos serán de exploración
\State \quad \textbf{[X4]} arranque en frío: sortea parámetros y calcula los soportes hasta \code{fecha_inicio}
\State \quad \textbf{[X4]} avanza vela a vela ejecutando los pasos A--F
\State \quad \textbf{[X4]} cada operación cerrada $\to$ una fila \code{'oc'} en \code{BTCUSD_store.csv}
\State \quad \textbf{[X4]} cada 4 velas con posiciones abiertas $\to$ una fila \code{'periodico'}
\State \quad \textbf{[X4]} cada 5 días simulados $\to$ parámetros nuevos y soportes recalculados
\State \quad \textbf{[X4]} llega al final de los datos (o quiebra); guarda checkpoint y termina
\State \textbf{[X5]} cuenta las filas \code{'oc'} del store
\If{$\ge 500$} \textbf{[X5]} entrena LightGBM y guarda modelo, columnas y métricas \EndIf
\State \textbf{[X5]} empieza el \textbf{ciclo 2}
\State \quad \textbf{[X4$\to$X5]} ahora sí hay modelo: 70\,\% de los tramos usarán inferencia
\State \quad \textbf{[X4]} en cada tramo de explotación construye el contexto \textit{al instante simulado} y pide parámetros a X5
\State \quad \textbf{[X4]} las filas nuevas del store ya reflejan decisiones informadas por el modelo
\State \dots\ hasta agotar \code{N_CICLOS_BT} o superar 5\,000 operaciones cerradas
\end{pseudo}

\subsection{Aislamiento: qué separa a qué}

Tres ejes de aislamiento operan simultáneamente y es fácil confundirlos:

\begin{enumerate}[leftmargin=1.6em,itemsep=3pt]
\item \textbf{Por activo.} \code{--activo BTCUSD} redirige las carpetas de trabajo de X4 a
\code{bt_BTCUSD}. Permite seis procesos simultáneos sin colisión. El store ya era por activo
por diseño.
\item \textbf{Por espacio de nombres.} La variable de entorno \code{X5_DEMO_SUFFIX} añade
\code{_demo} al store, al directorio del modelo, al historial de métricas y a las carpetas
de backtesting. Es lo que permite que el modo demostración no toque nada real.
\item \textbf{Por versión.} Las versiones V0, V1, \dots\ de X4 tienen cada una su
configuración congelada y su propia carpeta de resultados. Este eje es independiente de los
otros dos: el pipeline de X5 no usa versiones.
\end{enumerate}

\subsection{Supuestos y riesgos conocidos}

Vale la pena dejar por escrito lo que puede salir mal, porque nada de esto es un defecto que
se pueda «arreglar» sin más: son propiedades del enfoque.

\paragraph{El modelo aprende de datos que él mismo generó.} A partir del segundo ciclo, el
70\,\% de los tramos usa parámetros recomendados por el modelo. Si el modelo se equivoca de
forma sistemática, generará datos que confirmen su error. La exploración permanente del
30\,\% es el contrapeso, pero es un contrapeso, no una garantía.

\paragraph{El backtest no es la realidad.} Todo lo que el modelo aprende viene de una
simulación con supuestos optimistas: sin comisiones, sin deslizamiento, con intra-vela
sintética. Los parámetros óptimos \textit{para el simulador} pueden no serlo para el mercado.
La Fase 2 —realimentar con operaciones reales— existe precisamente para corregir eso.

\paragraph{Ocho parámetros y un objetivo ruidoso.} El retorno de una operación individual
tiene una componente de azar enorme. El modelo puede necesitar muchísimos datos para separar
la señal del ruido, y un $R^2$ negativo en las primeras iteraciones es lo esperable, no una
señal de fallo.

\paragraph{Contexto de cartera vacío en la inferencia en vivo.} En Fase 1, las columnas de
cartera se rellenan con ceros al inferir en vivo, mientras que en el store esas mismas
columnas traen valores reales. Es un desajuste entre entrenamiento e inferencia que solo se
resuelve cuando X1 provea el estado real de la cartera.

\paragraph{Las columnas OA se aproximan con las de OE.} Al inferir en vivo se copian los
valores del momento de crear la orden como si fueran los del momento de ejecutarla. Es una
aproximación razonable para órdenes que se ejecutan pronto y mala para las que tardan
semanas.

\subsection{Cómo saber si el bucle está funcionando}

Tres señales, en orden de disponibilidad:

\begin{enumerate}[leftmargin=1.6em,itemsep=2pt]
\item \code{python scripts/X5_macro_brain.py --status} muestra el recuento de operaciones
cerradas por activo y qué modelo corresponde. Si no crece entre recolecciones, X4 no está
cerrando posiciones y hay que mirar los parámetros o los datos antes que el modelo.
\item El historial en \code{resources/x5/Performance/\{ACTIVO\}_performance.json} muestra la
evolución de las métricas de validación entrenamiento tras entrenamiento. La señal buena es
que el error de validación baje mientras el store crece.
\item \code{equity_global.csv} de las carpetas \code{bt_\{ACTIVO\}} da la curva de capital de
la última pasada. Sirve para juzgar si las configuraciones que el modelo recomienda producen
pasadas mejores que las de exploración pura.
\end{enumerate}

\newpage

% ═══════════════════════════════════════════════════════════════════════════
\section{Anexos}
% ═══════════════════════════════════════════════════════════════════════════

\subsection{Glosario técnico}

\begin{longtable}{@{}p{4.3cm} p{\dimexpr\linewidth-4.8cm\relax}@{}}
\toprule
\textbf{Término} & \textbf{Definición} \\
\midrule
\endhead
Arranque en caliente & Iniciar un algoritmo iterativo desde una solución previa conocida en vez de un punto aleatorio. \\
Ascenso por gradiente & Optimizar moviéndose repetidamente en la dirección de máximo crecimiento de una función derivable. \\
Atención & Mecanismo por el que cada elemento de una secuencia pondera cuánto atender a los demás. Base del transformer. \\
Backtesting & Evaluar una estrategia sobre datos históricos, simulando que se opera en cada momento del pasado. \\
Escritura atómica & Escribir a un archivo temporal y renombrarlo sobre el destino, de modo que nunca quede a medio escribir. \\
FT-Transformer & Arquitectura transformer adaptada a datos tabulares: cada columna se convierte en un vector antes de aplicar atención. \\
Gradient boosting & Método que construye árboles de decisión en secuencia, cada uno corrigiendo el error de los anteriores. \\
Hueco (\textit{gap}) & Salto de precio entre el cierre de una vela y la apertura de la siguiente, sin cotizaciones intermedias. \\
LightGBM & Implementación eficiente de gradient boosting; estándar en datos tabulares. \\
Margen & Garantía inmovilizada al abrir una posición apalancada: valor de la posición dividido por el apalancamiento. \\
Nivel de margen & Patrimonio dividido por margen usado, en porcentaje. Indicador de supervivencia de la cuenta. \\
Modelo sustituto & Modelo barato que imita una función cara de evaluar, para poder optimizarla. \\
Normalización de capa & Reescalado a media cero y desviación uno dentro de cada vector; estabiliza redes profundas. \\
Optuna / TPE & Biblioteca de optimización de caja negra que concentra los intentos en las zonas prometedoras del espacio. \\
Orden límite de compra & Instrucción de comprar automáticamente si el precio baja hasta un nivel dado. \\
Parada temprana & Detener el entrenamiento cuando el error de validación deja de mejorar, para evitar sobreajuste. \\
Patrimonio (\textit{equity}) & Saldo realizado más el resultado flotante de las posiciones abiertas. \\
$R^2$ & Fracción de la variabilidad del objetivo explicada por el modelo. Negativo significa peor que predecir la media. \\
Sobreajuste & Memorizar el ruido del conjunto de entrenamiento, a costa del rendimiento con datos nuevos. \\
Trailing stop & Stop de venta cuyo nivel sube con el precio y nunca baja. \\
Vela (OHLCV) & Resumen de un intervalo de tiempo: apertura, máximo, mínimo, cierre y volumen. \\
\bottomrule
\end{longtable}

\subsection{Mapa de archivos generados}

\begin{longtable}{@{}p{6.6cm} p{\dimexpr\linewidth-7.1cm\relax}@{}}
\toprule
\textbf{Ruta} & \textbf{Contenido} \\
\midrule
\endhead
\code{resources/x4/version\{V\}/} & Configuración congelada y resultados de una versión de backtesting. \\
\code{\dots/resources_\{V\}/checkpoint.json} & Estado completo para reanudar. \\
\code{\dots/resources_\{V\}/trades.json} & Una entrada por operación cerrada. \\
\code{\dots/resources_\{V\}/events.json} & Bitácora granular de creaciones, cancelaciones, stops y cierres. \\
\code{\dots/resources_\{V\}/equity_global.csv} & Saldo, patrimonio y margen por vela. \\
\code{\dots/resources_\{V\}/equity_activos.csv} & Ganancia cerrada, flotante y total por vela y activo. \\
\code{\dots/conjuntos_N/\{A\}_\{N\}_bt.json} & Caché de soportes indexado por fecha. \\
\code{\dots/conjuntos_N/\{A\}_\{N\}_bt_delta.json} & Paso adaptativo del optimizador de soportes. \\
\code{resources/x5/\{A\}_store.csv} & Conjunto de datos de entrenamiento del activo. \\
\code{resources/x5/models/\{A\}/} & Modelos entrenados y orden de columnas. \\
\code{resources/x5/Performance/\{A\}_performance.json} & Historial de métricas por entrenamiento. \\
\code{resources/x5/bt_\{A\}/} & Carpeta de backtesting aislada del activo (modo X5). \\
\code{resources/x5/demo\_plots/\{A\}/} & Gráficos de cada búsqueda de soportes del recorrido guiado. \\
\code{config/active_parameters.json} & Recomendación vigente por activo, con su estado de modelo. \\
\bottomrule
\end{longtable}

\subsection{Índice de funciones}

Referencia rápida para saltar del documento al código. Los números de línea corresponden al
estado del repositorio en la fecha de este documento y son orientativos.

\subsubsection{\code{X4_backtester.py}}

\begin{longtable}{@{}p{5.2cm} r p{\dimexpr\linewidth-7.3cm\relax}@{}}
\toprule
\textbf{Función} & \textbf{Línea} & \textbf{Propósito} \\
\midrule
\endhead
\code{_cargar_config} & 36 & Carga la configuración congelada de una versión. \\
\code{_cargar_config_x5} & 47 & Carga \code{config_x5_default.py} para el modo X5. \\
\code{_cargar_datos_h1} & 68 & Lee, ordena y recorta las velas horarias por activo. \\
\code{_cargar_datos_m1} & 89 & Ídem para velas de un minuto; admite ausencia. \\
\code{_cargar_checkpoint} & 107 & Restaura el estado, reconvirtiendo claves de texto a número. \\
\code{_guardar_checkpoint} & 120 & Vuelca el estado a JSON. \\
\code{_worker_recalcular} & 144 & Adaptador que invoca \code{X0._procesar_valor_N} en un proceso aparte. \\
\code{_recalcular_soportes} & 228 & Orquesta el recálculo paralelo y recarga los soportes del caché. \\
\code{_calcular_estado_cuenta} & 267 & Saldo, patrimonio, margen usado y libre, nivel de margen. \\
\code{_id_orden} & 310 & Identificador reconstruible de una operación. \\
\code{_registrar_trade} & 320 & Arma la entrada de \code{trades.json}. \\
\code{_paso_A} & 361 & Cancela órdenes en espera sin soporte vigente. \\
\code{_paso_B} & 380 & Ejecuta órdenes alcanzadas; trata huecos y control de margen. \\
\code{_paso_C} & 465 & Activa y arrastra el stop protector. \\
\code{_paso_D} & 495 & Cierra por pérdida máxima. \\
\code{_paso_E} & 533 & Cierra por stop tocado. \\
\code{_paso_F} & 570 & Crea órdenes en espera sobre soportes libres. \\
\code{_necesita_intravela} & 611 & Decide si la vela requiere resolución minuto a minuto. \\
\code{_escalar_bloque_m1} & 629 & Deforma un tramo M1 real para encajarlo en la vela horaria. \\
\code{_simular_intravela} & 655 & Repite los pasos B--F sobre 60 velas de un minuto. \\
\code{_procesar_candle} & 675 & Aplica los pasos A--F a una vela de un activo. \\
\code{_append_equity} & 708 & Escribe las dos curvas de patrimonio. \\
\code{_flush_json_list} & 745 & Volcado atómico con reintentos de operaciones y eventos. \\
\code{_features_temporales_ts} & 777 & Features de calendario, incluidos festivos bursátiles. \\
\code{_compute_x5_snapshot} & 828 & Fotograma de indicadores X3 y score X2 al instante simulado. \\
\code{_contexto_portfolio_x5} & 840 & Estado de la cartera del activo como features. \\
\code{_construir_fila_oc} & 862 & Convierte una operación cerrada en fila de entrenamiento. \\
\code{_construir_fila_periodica} & 902 & Fila de seguimiento con posiciones abiertas. \\
\code{_append_x5_store} & 932 & Añade una fila al store, con defensa contra archivos truncados. \\
\code{ejecutar_backtest} & 956 & Bucle principal de la simulación. \\
\code{_params_explore_activo} & 1249 & Sorteo de parámetros dentro de los rangos. \\
\code{_contexto_inferencia_x5} & 1273 & Contexto al instante simulado para pedir inferencia. \\
\code{_seleccionar_params_x5} & 1299 & Decide explorar o explotar por activo. \\
\code{_aplicar_params_x5} & 1352 & Inyecta los parámetros del tramo en la configuración viva. \\
\code{ejecutar_x5_ciclo} & 1372 & Prepara el aislamiento por activo y carga el modelo del ciclo. \\
\bottomrule
\end{longtable}

\subsubsection{\code{X5_macro_brain.py}}

\begin{longtable}{@{}p{5.6cm} r p{\dimexpr\linewidth-7.7cm\relax}@{}}
\toprule
\textbf{Función} & \textbf{Línea} & \textbf{Propósito} \\
\midrule
\endhead
\code{_cargar_store} & 128 & Lee el CSV del activo descartando filas malformadas. \\
\code{_n_oc} & 138 & Cuenta filas de tipo operación cerrada. \\
\code{_seleccionar_tipo} & 144 & Elige \code{untrained}, \code{lgbm} o \code{ftt} según el volumen. \\
\code{_params_baseline} & 153 & Valores de \code{config.py}, usados como reserva. \\
\code{_feature_cols_de_store} & 168 & Fija el orden canónico de columnas y excluye las de cierre. \\
\code{_preparar_features} & 192 & Construye $X$, $y$ y los pesos temporales. \\
\code{_leer_contexto_actual} & 302 & Contexto de mercado en vivo. \\
\code{_construir_vector} & 336 & Combina contexto fijo y parámetros variables en el orden correcto. \\
\code{_calcular_metricas} & 346 & MAE, MSE, MAPE y $R^2$. \\
\code{_guardar_performance} & 365 & Añade una entrada al historial de métricas. \\
\code{_entrenar_lgbm} & 396 & Entrena los tres modelos de árboles con división temporal. \\
\code{FTTransformer} & 520 & Red con tronco compartido y tres cabezas de salida. \\
\code{_entrenar_ftt} & 558 & Entrena la red, normalizando con estadísticos del tramo de entrenamiento. \\
\code{_aplicar_airbag} & 679 & Fuerza exposición mínima tras una caída brusca. \\
\code{_inferir_lgbm} & 718 & Búsqueda bayesiana de parámetros con Optuna. \\
\code{_inferir_ftt} & 771 & Ascenso por gradiente sobre las entradas, con reinicios. \\
\code{_entrenar} & 876 & Elige y entrena el modelo que corresponda. \\
\code{cargar_modelo_para_activo} & 893 & Carga el modelo una sola vez, para reuso desde X4. \\
\code{inferir_con_contexto} & 910 & Inferencia con contexto inyectado, sin airbag. \\
\code{_escribir_active_parameters} & 948 & Mezcla y escritura atómica del archivo de parámetros. \\
\code{_mostrar_status} & 976 & Diagnóstico por activo. \\
\code{_worker_recolectar_activo} & 1018 & Bucle de ciclos y reentrenamiento de un activo. \\
\code{_recolectar} & 1090 & Orquestador paralelo de la recolección. \\
\code{_borrar_checkpoints_activo} & 1196 & Borrado completo del activo en su espacio de nombres. \\
\code{_recolectar_guiado} & 1248 & Recorrido secuencial narrado, en modo demostración u oficial. \\
\code{main} & 1365 & Despacho de los tres modos de la línea de comandos. \\
\bottomrule
\end{longtable}

\vspace{1cm}
\begin{center}
\small\color{cgris}
Documento generado a partir del código en \code{scripts/X4_backtester.py},
\code{scripts/X5_macro_brain.py}, \code{scripts/config_x5_default.py} y
\code{scripts/config.py}.\\
Fuente \LaTeX: \code{docs/plans/algoritmos.tex}. Compilar con \code{tectonic algoritmos.tex}.
\end{center}

\end{document}
